\chapter{The Spectral Flow $N_{\mathrm{BC}}\rightsquigarrow H_{\mathrm{HP}}$: Irreducible Obstruction at the Archimedean Place}
\label{v4-arith:w9-r4:hpar-spectral-flow}

\emph{Chapter~\ref{v4-arith:w5-b:hp-operator} delivered the Bost--Connes
log-modular operator $N_{\mathrm{BC}}=-\tfrac12\log\Delta_{1}$ with
unconditional self-adjointness and pure point spectrum
$\{\log m\mid m\in\mathbb{N}^{\times}\}$. The target operator is the
Hilbert--P\'olya ordinate $H_{\mathrm{HP}}$ with spectrum equal to the
imaginary parts $\{\gamma_{\rho}\}$ of non-trivial zeros of
$\xi_{\mathbb{R}}$. The intertwining operator $A$ taking
$N_{\mathrm{BC}}$ to $H_{\mathrm{HP}}$ is the classical
Hilbert--P\'olya twist. One hundred years of effort have not produced
it. This chapter states the spectral-flow problem rigorously, names
every obstruction in sequence, constructs the pieces that are
unconditionally available, isolates the residual as a specific
operator-theoretic statement equivalent to the Connes archimedean
absorption spectrum, and places it in the triple-identification
framework of Chapters~\ref{v4-arith:w8-a:deninger-chiral-homology},
\ref{v4-arith:w7-j:ATP-adelic-scattering}, and
\ref{v4-arith:3d-acs-E1bar}. No closure of H-P$^{\mathrm{Ar}}$ is
claimed. The chapter is an honest location document.}

\section{Setup: $N_{\mathrm{BC}}$ and $\Theta_{\xi}$ as Target Operators}
\label{v4-arith:w9-r4:setup}

We fix the two endpoints of the spectral flow and the Hilbert spaces
that carry them.

\subsection{Source: The Bost--Connes Log-Prime Operator}
\label{v4-arith:w9-r4:setup:source}

\begin{definition}[Bost--Connes Hilbert space and log-modular operator]
\label{v4-arith:w9-r4:def:bc-data}
\ClaimStatusDefinitional. Let
$(\mathcal{H}_{1},\pi_{1},\Omega_{1})$ denote the GNS triple of the
Bost--Connes KMS$_{1}$ state $\varphi_{1}$ on the adelic
$C^{*}$-algebra $C^{*}_{\mathbb{Q}}$
(\Cref{v4-arith:w5-b:def:gns},
\Cref{v4-arith:w5-b:lem:gns-restricted-tensor}). Let
$\mathcal{H}_{\mathrm{BC}}$ denote the closed subspace of
$\mathcal{H}_{1}$ spanned by the zero-mode lattice
$\{\delta_{m}\otimes\delta_{0}\mid m\in\mathbb{N}^{\times}\}$
restricted to the spherical character sector $\gamma=0$. Let
$N_{\mathrm{BC}}=-\tfrac12\log\Delta_{1}|_{\mathcal{H}_{\mathrm{BC}}}$
be the restriction of the log-modular operator; it is self-adjoint,
non-negative, with pure point spectrum $\{\log m\}_{m\geq 1}$ and
simple multiplicity on the spherical sector
(\Cref{v4-arith:w5-b:thm:N-BC-self-adjoint}).
\end{definition}

\begin{remark}[intrinsic structure of $\mathcal{H}_{\mathrm{BC}}$]
\label{v4-arith:w9-r4:rem:bc-structure}
As a Hilbert space, $\mathcal{H}_{\mathrm{BC}}$ is canonically
isometric to $\ell^{2}(\mathbb{N}^{\times},dm/m)$ with
$\delta_{m}\otimes\delta_{0}\leftrightarrow\sqrt{m}\,e_{m}$. Under
$m=\prod_{p}p^{v_{p}(m)}$ the lattice factors as
$\bigotimes'_{p}\ell^{2}(p^{\mathbb{N}_{\geq 0}},p^{-k}\cdot dk)$
restricted with respect to the local unit vector $e_{0}$ at each
place, with $N_{\mathrm{BC}}$ acting as
$\bigoplus_{p}k\log p$ under the Mellin isomorphism
$\ell^{2}(p^{\mathbb{N}_{\geq 0}},p^{-k}\,dk)\cong L^{2}(\mathbb{N},
p^{-k}\,dk)$. This is the Julia primon-gas presentation of
Chapter~\ref{v4-arith:connes-soibelman}, realised as a
Hilbert space of primon occupation states.
\end{remark}

\subsection{Target: The Hypothetical Zeta-Zero Operator}
\label{v4-arith:w9-r4:setup:target}

\begin{definition}[target Hilbert space and zero operator]
\label{v4-arith:w9-r4:def:target-data}
\ClaimStatusDefinitional. Let $\mathcal{Z}$ denote the multiset of
non-trivial zeros $\rho=1/2+i\gamma_{\rho}$ of $\xi_{\mathbb{R}}$ with
multiplicities $m_{\rho}$. The \emph{formal target Hilbert space} is
\[
\mathcal{H}_{\xi}^{\mathrm{formal}}\;:=\;\bigoplus_{\rho\in\mathcal{Z}}\mathbb{C}^{m_{\rho}},
\]
with inner product making the one-dimensional eigenspaces
orthonormal. The \emph{formal zeta-zero ordinate operator} is the
diagonal self-adjoint operator
\[
H_{\mathrm{HP}}^{\mathrm{formal}}\;:\;(\xi_{\rho})_{\rho}\;\mapsto\;(\gamma_{\rho}\xi_{\rho})_{\rho}
\]
on the finite-support core; its spectrum is
$\{\gamma_{\rho}\}$ with multiplicities. The \emph{formal zero
operator} is $\Theta_{\xi}^{\mathrm{formal}}=\tfrac12\mathrm{id}+iH_{\mathrm{HP}}^{\mathrm{formal}}$.
\end{definition}

\begin{remark}[why ``formal'']
\label{v4-arith:w9-r4:rem:why-formal}
The triple $(\mathcal{H}_{\xi}^{\mathrm{formal}},H_{\mathrm{HP}}^{\mathrm{formal}},\Theta_{\xi}^{\mathrm{formal}})$
exists trivially as a diagonal construction. It carries no
intrinsic arithmetic structure: the enumeration of
$\{\gamma_{\rho}\}$ is the classical index-set of non-trivial zeros,
and the diagonal operator realises its spectrum by fiat. The content
of the Hilbert--P\'olya programme is not the existence of
$\mathcal{H}_{\xi}^{\mathrm{formal}}$; it is the existence of a
Hilbert space \emph{with independent structure} whose spectral data
happen to agree with this multiset.
\end{remark}

\begin{definition}[non-formal target]
\label{v4-arith:w9-r4:def:non-formal-target}
\ClaimStatusDefinitional. A \emph{non-formal target} is a pair
$(\mathcal{H}_{\xi},H_{\mathrm{HP}})$ where $\mathcal{H}_{\xi}$ is a
Hilbert space carrying structure inherited from the Bost--Connes
framework (restricted-tensor adelic factorisation, Petersson inner
product on cuspidal cores, GNS modular data, or another
independently specified structure) and $H_{\mathrm{HP}}$ is a
self-adjoint operator on $\mathcal{H}_{\xi}$ with spectrum
$\{\gamma_{\rho}\}$ whose spectral resolution is canonical with
respect to the inherited structure.
\end{definition}

\begin{remark}[scope of ``canonical'']
\label{v4-arith:w9-r4:rem:canonical-scope}
``Canonical'' means: the spectral resolution of $H_{\mathrm{HP}}$ is
determined by the inherited structure up to a unitary equivalence
whose parameters are themselves determined by that structure. A
diagonal construction with eigenvalues labelled by an arbitrary
enumeration is not canonical; a diagonal construction whose
eigenvectors are explicit functions of arithmetic data (adelic
characters, automorphic forms, spectral projectors of a globally
defined operator) is.
\end{remark}

\subsection{Spectral Flow as Intertwining Operator}
\label{v4-arith:w9-r4:setup:flow}

\begin{definition}[the spectral flow problem]
\label{v4-arith:w9-r4:def:spectral-flow}
\ClaimStatusDefinitional. The \emph{spectral flow problem} from
$N_{\mathrm{BC}}$ to $H_{\mathrm{HP}}$ is the specification of:
\begin{enumerate}
\item a Hilbert space $\mathcal{H}_{\xi}$ with a non-formal structure
in the sense of \Cref{v4-arith:w9-r4:def:non-formal-target};
\item a densely defined linear operator
$A:\mathcal{D}(A)\subset\mathcal{H}_{\mathrm{BC}}\to\mathcal{H}_{\xi}$
with dense range;
\item a self-adjoint operator $H_{\mathrm{HP}}$ on $\mathcal{H}_{\xi}$
with domain $\mathcal{D}(H_{\mathrm{HP}})\supset\mathrm{Ran}(A)$;
\item the intertwining identity
\begin{equation}
\label{v4-arith:w9-r4:eq:intertwine}
A\,N_{\mathrm{BC}}\,A^{-1}\;=\;H_{\mathrm{HP}}\qquad\text{on }\mathrm{Ran}(A)\cap\mathcal{D}(H_{\mathrm{HP}}),
\end{equation}
interpreted in the closed-operator sense after specifying domain
compatibility;
\item the spectral identity
\begin{equation}
\label{v4-arith:w9-r4:eq:spec-id}
\mathrm{Spec}(H_{\mathrm{HP}})\;=\;\{\gamma_{\rho}\mid\rho\in\mathcal{Z}\}
\end{equation}
with multiplicities matching.
\end{enumerate}
A \emph{solution} is a tuple $(\mathcal{H}_{\xi},A,H_{\mathrm{HP}})$
satisfying all five items. The \emph{Hilbert--P\'olya twist} is a
solution.
\end{definition}

\begin{remark}[unitarity of $A$]
\label{v4-arith:w9-r4:rem:unitarity-of-A}
\Cref{v4-arith:w9-r4:def:spectral-flow} does \emph{not} require $A$
to be bounded or unitary. The closed-operator intertwining
\eqref{v4-arith:w9-r4:eq:intertwine} admits unbounded $A$; in the
regular case one asks for a polar decomposition $A=U|A|$ with $U$
partial isometry and $|A|$ self-adjoint non-negative, on which the
intertwining reads
$U|A|N_{\mathrm{BC}}|A|^{-1}U^{*}=H_{\mathrm{HP}}$. Under
essential self-adjointness of $H_{\mathrm{HP}}$ and
$N_{\mathrm{BC}}$ the intertwining is equivalent to
$U\,\exp(it N_{\mathrm{BC}})\,U^{*}=\exp(itH_{\mathrm{HP}})$ on a
common core, i.e.\ a unitary intertwining of the one-parameter
groups. We use the closed-operator form throughout; when the
literature writes ``$A$ is a unitary twist'' it is the polar $U$ in
this decomposition.
\end{remark}

\section{Prime-Zero Mangoldt-Hadamard Duality Recap}
\label{v4-arith:w9-r4:mangoldt-hadamard}

The formal bridge between the two spectra is the Riemann-von Mangoldt
explicit formula and its Hadamard inversion. This section records
the bridge at the level of distributions; it is not an operator
theorem.

\subsection{Hadamard Factorisation of $\xi_{\mathbb{R}}$}
\label{v4-arith:w9-r4:mh:hadamard}

\begin{proposition}[Hadamard factorisation]
\label{v4-arith:w9-r4:prop:hadamard}
\ClaimStatusProvedElsewhere \emph{(Hadamard 1893;
\Cref{v4-arith:w5-b:eq:hadamard})}. The completed Riemann zeta
$\xi_{\mathbb{R}}(s)=\tfrac12 s(s-1)\pi^{-s/2}\Gamma(s/2)\zeta(s)$ is
entire of order $1$, of genus $1$, and admits the factorisation
\begin{equation}
\label{v4-arith:w9-r4:eq:hadamard-factor}
\xi_{\mathbb{R}}(s)\;=\;\xi_{\mathbb{R}}(0)\,e^{Bs}\,\prod_{\rho}\Bigl(1-\frac{s}{\rho}\Bigr)e^{s/\rho},
\end{equation}
with $\rho$ ranging over non-trivial zeros with multiplicity and
$B=-\log 2\pi+1+\gamma/2$ (Edwards 1974 \S2.7) the Hadamard constant.
\end{proposition}

\subsection{Riemann-von Mangoldt Explicit Formula}
\label{v4-arith:w9-r4:mh:rvm}

\begin{proposition}[Riemann-von Mangoldt explicit formula]
\label{v4-arith:w9-r4:prop:rvm}
\ClaimStatusProvedElsewhere \emph{(von Mangoldt 1895; Edwards 1974 \S3)}.
For $x>1$ non-integer,
\begin{equation}
\label{v4-arith:w9-r4:eq:rvm}
\sum_{n\leq x}\Lambda(n)\;=\;x\;-\;\sum_{\rho}\frac{x^{\rho}}{\rho}\;-\;\frac{\zeta'(0)}{\zeta(0)}\;-\;\tfrac12\log(1-x^{-2}),
\end{equation}
where $\Lambda(n)=\log p$ if $n=p^{k}$ and $0$ otherwise.
\end{proposition}

\subsection{Mellin-Dual Measures}
\label{v4-arith:w9-r4:mh:mellin-dual}

\begin{definition}[prime and zero counting measures]
\label{v4-arith:w9-r4:def:counting-measures}
\ClaimStatusDefinitional. The \emph{prime-counting distribution} is
\begin{equation}
\label{v4-arith:w9-r4:eq:mu-prime}
\mu_{\mathrm{prime}}\;:=\;\sum_{m\in\mathbb{N}^{\times}}\Lambda(m)\,\delta_{\log m}\;=\;\sum_{p\in\mathbb{P}}\sum_{k\geq 1}(\log p)\,\delta_{k\log p}
\end{equation}
on $\mathbb{R}_{\geq 0}$. The \emph{zero-counting distribution} is
\begin{equation}
\label{v4-arith:w9-r4:eq:mu-zero}
\mu_{\mathrm{zero}}\;:=\;\sum_{\rho\in\mathcal{Z}}\delta_{\gamma_{\rho}}
\end{equation}
on $\mathbb{R}$.
\end{definition}

\begin{proposition}[Mellin-dual form of the explicit formula]
\label{v4-arith:w9-r4:prop:mellin-dual}
\ClaimStatusProvedElsewhere \emph{(Weil 1952; Patterson 1988 Ch.~3)}.
For every even Schwartz function $h\in\mathcal{S}(\mathbb{R})$ with
Fourier transform $\widehat{h}$ of rapid decay,
\begin{equation}
\label{v4-arith:w9-r4:eq:mellin-dual}
\int_{\mathbb{R}}h\,d\mu_{\mathrm{zero}}\;=\;\int_{\mathbb{R}}h\,dW_{\infty}\;-\;2\sum_{m\geq 2}\frac{\Lambda(m)}{\sqrt{m}}\widehat{h}(\log m),
\end{equation}
with $W_{\infty}$ the archimedean Weil density of
\eqref{v4-arith:tfe:eq:W-arch}. The pairing is Mellin-dual: under
Mellin transform, the prime-side distribution
$\sum_{m}\Lambda(m)m^{-1/2}\delta_{\log m}$ is dual to the zero-side
distribution $\mu_{\mathrm{zero}}$ via the Fourier-Mellin pairing on
$\mathbb{R}$.
\end{proposition}

\begin{remark}[the duality is distributional, not operatorial]
\label{v4-arith:w9-r4:rem:duality-distributional}
The Mellin-Hadamard duality of
\Cref{v4-arith:w9-r4:prop:mellin-dual} is a statement about tempered
distributions on $\mathbb{R}$: both sides are linear functionals on
$\mathcal{S}(\mathbb{R})$, and the equality holds distributionally. It
is not an operator identity. In particular, the duality does not by
itself produce a linear map from one Hilbert space to another; it
produces a pairing between two distributions that, when integrated
against test functions, agree up to the archimedean correction.
This is the precise sense in which the classical Hilbert--P\'olya
philosophy ``interprets'' the duality as an operator-level
intertwining, and the reason the interpretation has not been made
rigorous: no canonical passage from distributions on $\mathbb{R}$ to
operators on a Hilbert space is known that realises both sides as
spectral projectors of a single self-adjoint operator.
\end{remark}

\section{Formal Spectral Flow $A$ as Mellin-Inversion Symbol}
\label{v4-arith:w9-r4:formal-flow}

We write down the candidate spectral flow operator at the symbolic
level, then identify what prevents it from being a genuine operator
between Hilbert spaces.

\subsection{The Mellin-Hadamard Inversion Symbol}
\label{v4-arith:w9-r4:formal-flow:symbol}

\begin{definition}[formal Mellin-Hadamard inversion]
\label{v4-arith:w9-r4:def:formal-MH}
\ClaimStatusDefinitional. The \emph{formal Mellin-Hadamard inversion}
is the composite symbol
\begin{equation}
\label{v4-arith:w9-r4:eq:formal-MH}
A^{\mathrm{sym}}\;:=\;\mathcal{I}_{\mathrm{Had}}\circ\mathcal{M}
\end{equation}
where $\mathcal{M}$ is the Mellin transform
$(\mathcal{M}f)(s)=\int_{0}^{\infty}f(x)x^{s-1}dx$ on a suitable
function class and $\mathcal{I}_{\mathrm{Had}}$ is the
Hadamard-inversion formal symbol
\[
(\mathcal{I}_{\mathrm{Had}}F)(\gamma_{\rho})\;:=\;\operatorname*{Res}_{s=\rho}\frac{F(s)}{\xi_{\mathbb{R}}(s)}
\]
reading the residue of $F/\xi_{\mathbb{R}}$ at each non-trivial zero
and writing it into the $\rho$-component of a formal sequence
indexed by $\mathcal{Z}$.
\end{definition}

\begin{remark}[$A^{\mathrm{sym}}$ is not an operator]
\label{v4-arith:w9-r4:rem:Asym-not-operator}
$A^{\mathrm{sym}}$ is a symbolic map from functions on
$\mathbb{R}_{>0}$ to sequences indexed by $\mathcal{Z}$. For it to
be an operator between Hilbert spaces one needs:
\begin{itemize}
\item a Hilbert space domain $\mathcal{H}_{\mathrm{BC}}\subset L^{2}$-functions
on $\mathbb{R}_{>0}$ where $\mathcal{M}$ is densely defined;
\item a Hilbert space codomain $\mathcal{H}_{\xi}\subset\ell^{2}(\mathcal{Z})$
where $\mathcal{I}_{\mathrm{Had}}$ lands;
\item a compatibility condition making $A^{\mathrm{sym}}$ closable.
\end{itemize}
Each of these is a non-trivial construction.
\end{remark}

\subsection{Formal Mellin--Hadamard Obstruction}
\label{v4-arith:w9-r4:formal-flow:formal-intertwine}

\begin{lemma}[formal Mellin--Hadamard obstruction]
\label{v4-arith:w9-r4:lem:formal-intertwine}
\ClaimStatusProvedHere. Suppose $f\in C_{c}^{\infty}(\mathbb{R}_{>0})$
has Mellin transform $F(s)=(\mathcal{M}f)(s)$ entire, and suppose
$\rho\in\mathcal{Z}$ is a simple zero of \(\xi_{\mathbb R}\). Then
\begin{equation}
\label{v4-arith:w9-r4:eq:formal-intertwine}
A^{\mathrm{sym}}\bigl(N_{\mathrm{BC}}f\bigr)(\gamma_{\rho})
\;=\;
-\,\frac{F'(\rho)}{\xi_{\mathbb R}'(\rho)},
\qquad
A^{\mathrm{sym}}(f)(\gamma_{\rho})
\;=\;
\frac{F(\rho)}{\xi_{\mathbb R}'(\rho)}.
\end{equation}
Consequently the naive formal intertwining relation
\[
A^{\mathrm{sym}}N_{\mathrm{BC}}
\stackrel{?}{=}
M_{\gamma}A^{\mathrm{sym}}
\]
holds at \(\rho\) only under the additional interpolation condition
\[
F'(\rho)=-\gamma_{\rho}F(\rho).
\]
This condition is not a consequence of RH.
\end{lemma}

\begin{proof}
Apply the Mellin transform: $N_{\mathrm{BC}}$ acts on
$\mathbb{R}_{>0}$-functions as
$(N_{\mathrm{BC}}f)(x)=-\log x\cdot f(x)$ in the logarithmic scale
(on the zero-mode lattice
\Cref{v4-arith:w9-r4:rem:bc-structure} the eigenvalue $\log m$ on
$e_{m}$ is the image under Mellin of the scaling by $m$). Mellin
transform gives
$(\mathcal{M}(N_{\mathrm{BC}}f))(s)=-F'(s)$, where $F'$ is the
derivative in the spectral variable. Since \(\rho\) is a simple zero,
\[
\operatorname*{Res}_{s=\rho}\frac{F(s)}{\xi_{\mathbb R}(s)}
=\frac{F(\rho)}{\xi_{\mathbb R}'(\rho)},
\qquad
\operatorname*{Res}_{s=\rho}\frac{-F'(s)}{\xi_{\mathbb R}(s)}
=-\frac{F'(\rho)}{\xi_{\mathbb R}'(\rho)}.
\]
Multiplication by the zero ordinate would instead give
\(\gamma_{\rho}F(\rho)/\xi_{\mathbb R}'(\rho)\). Equality is exactly
the displayed first-order interpolation condition.
\end{proof}

\begin{remark}[RH is not enough for formal intertwining]
\label{v4-arith:w9-r4:rem:formal-assumes-RH}
\Cref{v4-arith:w9-r4:lem:formal-intertwine} shows that the Mellin
generator differentiates the numerator before Hadamard residue
extraction. RH would place \(\rho\) on the critical line, but it would
not imply \(F'(\rho)=-\gamma_{\rho}F(\rho)\) for the Mellin transforms
in the source class. Thus the naive symbol
\(A^{\mathrm{sym}}\) is not a spectral-flow operator. A genuine
operator construction must supply a new domain theorem or boundary
condition that converts derivative residue data into zero-ordinate
multiplication without inserting that conclusion by definition.
\end{remark}

\section{Obstacle 1: Hilbert-Space Completions on Both Sides}
\label{v4-arith:w9-r4:obstacle-hilbert}

The first obstruction is the choice of Hilbert-space completions on
which the Mellin-Hadamard symbol lands as a closable operator. Both
sides have canonical candidates; the candidates are mutually
incompatible.

\subsection{Source-Side Completion}
\label{v4-arith:w9-r4:obstacle-hilbert:source}

\begin{definition}[source-side candidate Hilbert space]
\label{v4-arith:w9-r4:def:source-completion}
\ClaimStatusDefinitional. The source-side completion is
$\mathcal{H}_{\mathrm{BC}}=\ell^{2}(\mathbb{N}^{\times},dm/m)$ of
\Cref{v4-arith:w9-r4:rem:bc-structure}. $N_{\mathrm{BC}}$ acts as
multiplication by $\log m$ and is essentially self-adjoint on the
finite-support core
(\Cref{v4-arith:w5-b:thm:N-BC-self-adjoint}).
\end{definition}

\begin{lemma}[no Mellin transform on $\mathcal{H}_{\mathrm{BC}}$]
\label{v4-arith:w9-r4:lem:no-mellin-on-HBC}
\ClaimStatusProvedHere. The Mellin transform
$\mathcal{M}:L^{2}(\mathbb{R}_{>0},d^{\times}x)\to L^{2}(\mathbb{R},dt)$
(with $s=1/2+it$) is a classical unitary isomorphism
(Mellin 1904; Titchmarsh 1937 Ch.~7). Its restriction to
$\mathcal{H}_{\mathrm{BC}}\subset L^{2}(\mathbb{R}_{>0},d^{\times}x)$,
identified via
$e_{m}\leftrightarrow\delta_{x=m}$, is a map into the space of
``purely atomic'' distributions, which does not land in
$L^{2}(\mathbb{R},dt)$: the Mellin transform of a Dirac delta at
$m$ is $m^{s-1}$, not an $L^{2}$-function.
\end{lemma}

\begin{proof}
$\mathcal{M}(\delta_{x=m})(s)=\int_{0}^{\infty}\delta_{x=m}x^{s-1}dx=m^{s-1}$,
which is a bounded function of $s$ on any vertical line
$\mathrm{Re}(s)=c$ but not integrable. Hence it is not an $L^{2}$-function.
The linear combinations $\sum_{m}c_{m}\delta_{x=m}$ with
$\sum_{m}|c_{m}|^{2}/m<\infty$ have Mellin transform
$\sum_{m}c_{m}m^{s-1}$, a Dirichlet series, which does not have an
$L^{2}(\mathbb{R},dt)$-norm in general.
\end{proof}

\begin{remark}[the lattice embedding is ill-suited to Mellin]
\label{v4-arith:w9-r4:rem:lattice-ill-suited}
The issue is that $\mathcal{H}_{\mathrm{BC}}$ embeds in $L^{2}(\mathbb{R}_{>0})$
only as a space of atomic measures, not as a space of locally
integrable functions. Mellin transform on atomic measures produces
Dirichlet series, and a Dirichlet series is a distribution, not an
$L^{2}$-function on the critical line. The Mellin-transform
intertwiner of \Cref{v4-arith:w9-r4:def:formal-MH} therefore cannot
be defined as a bounded operator from $\mathcal{H}_{\mathrm{BC}}$
into $L^{2}(\mathbb{R}_{1/2+it})$.
\end{remark}

\subsection{Target-Side Completion}
\label{v4-arith:w9-r4:obstacle-hilbert:target}

\begin{definition}[target-side candidates]
\label{v4-arith:w9-r4:def:target-completion}
\ClaimStatusDefinitional. Three candidate Hilbert spaces for
$\mathcal{H}_{\xi}$ in the non-formal sense:
\begin{enumerate}
\item $\mathcal{H}_{\xi}^{\mathrm{idele}}=L^{2}(C_{\mathbb{Q}},d^{\times}x)$
the idele class Hilbert space of
\Cref{v4-arith:w7-j:def:H-idele};
\item $\mathcal{H}_{\xi}^{\mathrm{adele}}=L^{2}(X_{\mathbb{Q}})$ the
adele-class space with non-commutative-geometric $L^{2}$-structure
via $C^{*}(\mathbb{A}/\mathbb{Q}^{\times})$
(\Cref{v4-arith:w7-j:def:connes-ac-space});
\item $\mathcal{H}_{\xi}^{\mathrm{Petersson}}$ the Petersson
completion of the sph-finite cuspidal core
(\Cref{v4-arith:thm:P3-resolution}).
\end{enumerate}
\end{definition}

\begin{lemma}[Mellin transform on $\mathcal{H}_{\xi}^{\mathrm{idele}}$]
\label{v4-arith:w9-r4:lem:mellin-on-idele}
\ClaimStatusProvedElsewhere \emph{(Meyer 2005 \S3.2;
\Cref{v4-arith:w7-j:lem:mellin-decomp})}. The Mellin transform on
$L^{2}(\mathbb{R}_{>0},d^{\times}t)$ into
$L^{2}(\mathbb{R}_{1/2+it},dt)$ is unitary; it extends to the
idele class Hilbert space $\mathcal{H}_{\xi}^{\mathrm{idele}}$ by
tensoring with the identity on $L^{2}(C_{\mathbb{Q}}^{1})$.
\end{lemma}

\begin{remark}[the completion mismatch]
\label{v4-arith:w9-r4:rem:completion-mismatch}
\Cref{v4-arith:w9-r4:lem:no-mellin-on-HBC} shows Mellin fails on the
source $\mathcal{H}_{\mathrm{BC}}$, and
\Cref{v4-arith:w9-r4:lem:mellin-on-idele} shows Mellin is unitary on
the target $\mathcal{H}_{\xi}^{\mathrm{idele}}$. The two completions
are \emph{not} isometric: $\mathcal{H}_{\mathrm{BC}}$ is a space of
atomic measures on primes; $\mathcal{H}_{\xi}^{\mathrm{idele}}$ is
an $L^{2}$ of the continuous idele norm. The passage from one to
the other requires either (a) a Schwartz-class density argument
lifting atomic measures to a common completion where Mellin is
densely defined (Meyer 2005 \S3 route) or (b) a direct
factorisation homology argument making the zero-mode lattice
into a dense subspace of a factorisation algebra whose Mellin
transform lives naturally on $\mathbb{R}_{1/2+it}$
(Costello-Gwilliam-style route). Neither is delivered unconditionally
by the Bost--Connes framework alone.
\end{remark}

\section{Obstacle 2: Kernel of $A$ as Bounded/Unbounded Operator}
\label{v4-arith:w9-r4:obstacle-kernel}

Even accepting a compatible Hilbert-space completion, $A$ must have
an explicit kernel making it closable. The Hadamard inversion step is
the bottleneck.

\subsection{The Hadamard Inversion Kernel}
\label{v4-arith:w9-r4:obstacle-kernel:had}

\begin{definition}[formal Hadamard kernel]
\label{v4-arith:w9-r4:def:formal-had-kernel}
\ClaimStatusDefinitional. The formal kernel of
$\mathcal{I}_{\mathrm{Had}}:$ space of entire functions on
$\mathbb{C}$ $\to$ sequences on $\mathcal{Z}$ is
\[
K_{\mathrm{Had}}(\rho,s)\;:=\;\frac{1}{2\pi i}\cdot\frac{1}{\xi_{\mathbb{R}}(s)}\cdot\frac{1}{s-\rho}
\]
interpreted in the residue sense at $s=\rho$.
\end{definition}

\begin{lemma}[the Hadamard kernel is not a Hilbert-Schmidt kernel]
\label{v4-arith:w9-r4:lem:had-not-HS}
\ClaimStatusProvedHere. For any pair of Hilbert-space completions
$\mathcal{H}_{\mathrm{BC}}$ and $\mathcal{H}_{\xi}$ compatible with
Mellin inversion on one side of the spectral-flow obligation, the
Hadamard kernel $K_{\mathrm{Had}}$ is not in
$\mathrm{HS}(\mathcal{H}_{\xi},\mathcal{H}_{\mathrm{BC}})$: its
Hilbert-Schmidt norm
$\|K_{\mathrm{Had}}\|_{\mathrm{HS}}^{2}=\sum_{\rho}\int_{\mathrm{Re}(s)=1/2}|\xi_{\mathbb{R}}(s)|^{-2}|s-\rho|^{-2}ds$
diverges on the critical line at every $\rho$.
\end{lemma}

\begin{proof}
$\xi_{\mathbb{R}}$ has a zero at $\rho$, so $|\xi_{\mathbb{R}}(s)|^{-2}\sim C_{\rho}|s-\rho|^{-2m_{\rho}}$
near $s=\rho$ on the critical line. The integrand near $s=\rho$ is
therefore $|s-\rho|^{-2m_{\rho}-2}$, which is non-integrable on the
1-dimensional critical line for any $m_{\rho}\geq 1$. Summing over
$\rho$ only increases the divergence.
\end{proof}

\begin{remark}[unboundedness is intrinsic]
\label{v4-arith:w9-r4:rem:unboundedness-intrinsic}
\Cref{v4-arith:w9-r4:lem:had-not-HS} means that $A$ cannot be a
bounded operator under any reasonable pair of completions. The
spectral flow is therefore necessarily an \emph{unbounded} operator,
requiring the machinery of closable operators with dense domain. The
choice of domain becomes a substantive part of the problem, not a
technicality.
\end{remark}

\subsection{The Domain Problem}
\label{v4-arith:w9-r4:obstacle-kernel:domain}

\begin{proposition}[candidate domains for $A$]
\label{v4-arith:w9-r4:prop:candidate-domains}
\ClaimStatusProvedHere. Three candidate domains for $A$ in
$\mathcal{H}_{\mathrm{BC}}$:
\begin{enumerate}
\item $\mathcal{D}_{\mathrm{fin}}$: finite-support vectors
$\sum_{m\leq M}c_{m}e_{m}$ for $M$ finite. On these, $A$ can be
computed via residue extraction; the result is a finite sum of
residues, well-defined.
\item $\mathcal{D}_{\mathrm{Schwartz}}$: Schwartz-class-like vectors
$\sum_{m}c_{m}e_{m}$ with $|c_{m}|=O(m^{-N})$ for all $N$. On these,
the Dirichlet series $F(s)=\sum_{m}c_{m}m^{s-1}$ extends entirely
with rapid decay on vertical strips; $A$ is defined distributionally
by residue contour deformation.
\item $\mathcal{D}_{\mathrm{BR}}$: the Beurling-Rajagopalan domain
of functions whose Mellin transform has exponential-polynomial
growth on the critical strip bounded by Plancherel-P\'olya.
\end{enumerate}
On $\mathcal{D}_{\mathrm{fin}}$ and $\mathcal{D}_{\mathrm{Schwartz}}$,
$A$ is densely defined but closability is conditional on the decay
rate of $c_{m}$ being compatible with Hadamard inversion.
\end{proposition}

\begin{proof}
Item (1): a finite sum $\sum_{m\leq M}c_{m}e_{m}$ maps under $A$ to
$\sum_{\rho}r_{\rho}(f)\cdot e_{\rho}^{\mathrm{target}}$ where
$r_{\rho}(f)=\operatorname{Res}_{s=\rho}(\sum_{m\leq M}c_{m}m^{s-1})/\xi_{\mathbb{R}}(s)$
is a finite expression involving derivatives of $\xi_{\mathbb{R}}^{-1}$
at $\rho$. The sum over $\rho$ is formal; summability depends on
$\sum_{\rho}|r_{\rho}|^{2}<\infty$, which holds for Schwartz $c_{m}$
but not in general.
Item (2): Schwartz decay on $c_{m}$ gives entire $F(s)$ with rapid
decay on vertical strips; residue deformation converts the sum
$\sum_{\rho}r_{\rho}(f)$ into a contour integral over a vertical line
shifted to the left, plus boundary terms at $s=0,1$ and at trivial
zeros. The contour-integral part is a Fourier-Mellin integral bounded
by Plancherel; the boundary terms are finite.
Item (3): Beurling 1955 (Ann.~Math.~61) and Rajagopalan 1960 show
that the Beurling-Rajagopalan class is a natural setting for Mellin
inversion with controlled exponential growth, used in Nyman-Beurling
reformulations of RH; see Chapter~\ref{v4-arith:w7-b:chapter}
of the present volume for the chiral-algebra adaptation.
\end{proof}

\begin{remark}[no domain gives self-adjointness]
\label{v4-arith:w9-r4:rem:no-domain-gives-SA}
None of the three candidate domains of
\Cref{v4-arith:w9-r4:prop:candidate-domains} closes to a domain on
which $A^{*}A=|A|^{2}$ has a canonical square-root polar decomposition
with unitary part $U$ satisfying the intertwining
\eqref{v4-arith:w9-r4:eq:intertwine}. Essential self-adjointness of
$A^{*}A$ on any of these domains is an open question, and its
failure is the functional-analytic face of the Hilbert--P\'olya
obstruction.
\end{remark}

\section{Connes 1999 Partial Approach and Its Limitation}
\label{v4-arith:w9-r4:connes-1999}

Connes' adele class space programme is the canonical partial
resolution of the spectral-flow problem. We record what it
constructs unconditionally and where it cannot close.

\subsection{The Connes Construction}
\label{v4-arith:w9-r4:connes-1999:construction}

\begin{definition}[Connes adele class data]
\label{v4-arith:w9-r4:def:connes-data}
\ClaimStatusDefinitional \emph{(Connes 1999 arXiv:math/9811068 \S IV)}.
The \emph{Connes adele class Hilbert space} is
$\mathcal{H}_{\mathrm{Connes}}=L^{2}(X_{\mathbb{Q}},dx)$ with
$X_{\mathbb{Q}}=\mathbb{A}/\mathbb{Q}^{\times}$ the adele class space
and $dx$ the Haar measure inherited from $\mathbb{A}$ modulo
$\mathbb{Q}^{\times}$; the quotient is non-Hausdorff, and the
$L^{2}$-structure is realised via the crossed product
$C^{*}(\mathbb{A}/\mathbb{Q}^{\times})=C_{0}(\mathbb{A})\rtimes\mathbb{Q}^{\times}$.
The \emph{Connes scaling generator} is
$D_{\mathrm{Connes}}=-i\partial_{\log\lambda}U(\lambda)$ at
$\lambda=1$, where $U(\lambda)f(x)=|\lambda|^{-1/2}f(\lambda^{-1}x)$.
On the cuspidal subspace
$\mathcal{E}\subset\mathcal{H}_{\mathrm{Connes}}$ (orthogonal to the
Tate-regularised constants), $D_{\mathrm{Connes}}$ is conjectured to
have pure point spectrum equal to the imaginary parts of non-trivial
zeros.
\end{definition}

\begin{theorem}[Connes trace formula equivalence]
\label{v4-arith:w9-r4:thm:connes-trace-eq}
\ClaimStatusProvedElsewhere \emph{(Connes 1999 \S V Thm.~1;
\Cref{v4-arith:tfe:thm:connes-meyer})}. For
$h\in\mathcal{H}_{\mathrm{test}}$ (the even-Schwartz test class on
$\mathbb{R}$), the regularised trace
$\mathrm{Tr}_{\mathrm{reg}}(h(D_{\mathrm{Connes}})|_{\mathcal{E}})$
equals the Weil distribution side, conditional on the
absorption-spectrum conjecture. Equivalently, the regularised trace
equals $\sum_{\rho}h(\gamma_{\rho})$ if and only if the non-trivial
zeros lie on the critical line.
\end{theorem}

\subsection{The Limitation: Absorption Spectrum as RH}
\label{v4-arith:w9-r4:connes-1999:limitation}

\begin{theorem}[Connes absorption-spectrum conjecture]
\label{v4-arith:w9-r4:thm:connes-absorption}
\ClaimStatusConjectured \emph{(Connes 1999 \S IV Thm.~1 conditional
form)}. The scaling generator $D_{\mathrm{Connes}}$ on
$\mathcal{E}\subset\mathcal{H}_{\mathrm{Connes}}$ has absorption
spectrum equal to $\{\gamma_{\rho}\mid\rho\in\mathcal{Z}\}$ with
multiplicities matching. Equivalently, the non-trivial zeros lie on
the critical line.
\end{theorem}

\begin{remark}[Connes does not construct $A$ unconditionally]
\label{v4-arith:w9-r4:rem:connes-not-uncond}
Connes 1999 does not construct the intertwining $A$
unconditionally. The Connes construction supplies a candidate
Hilbert space $\mathcal{H}_{\mathrm{Connes}}$ and a candidate
self-adjoint operator $D_{\mathrm{Connes}}$, but the identification
of $D_{\mathrm{Connes}}$'s absorption spectrum with the zero set is
equivalent to RH. In the spectral-flow vocabulary, Connes supplies
the triple $(\mathcal{H}_{\xi},A,H_{\mathrm{HP}})=(\mathcal{H}_{\mathrm{Connes}},\mathrm{id},D_{\mathrm{Connes}})$
with identity intertwiner, and shifts the question from
``construct $A$'' to ``prove $D_{\mathrm{Connes}}$ has the right
absorption spectrum''. The two questions have the same strength.
\end{remark}

\begin{proposition}[Connes equivalence to H-P$^{\mathrm{Ar}}$]
\label{v4-arith:w9-r4:prop:connes-eq-HPAr}
\ClaimStatusProvedElsewhere \emph{(Connes 1999 \S VII;
\Cref{v4-arith:hpvbconv:cor:HPAr-connes-equivalence} as recorded in
Ch.~\ref{v4-arith:HPAr-vbconverse})}. The Connes absorption-spectrum
conjecture
(\Cref{v4-arith:w9-r4:thm:connes-absorption}) is equivalent to the
Hilbert--P\'olya hypothesis
H-P$^{\mathrm{Ar}}$ (\Cref{v4-arith:rh:hyp:HPAr}) of the Vol~IV
arithmetic branch.
\end{proposition}

\begin{remark}[the adele-class route transfers the problem]
\label{v4-arith:w9-r4:rem:adele-transfers}
The Connes route transfers the spectral-flow problem from the
Bost--Connes GNS side to the adele-class side. The adele-class
side has an \emph{intrinsic} candidate for the Hilbert space
($L^{2}$ of the non-Hausdorff quotient) and an \emph{intrinsic}
candidate for the zero operator (scaling generator on cuspidal
subspace), but the intrinsic candidates are equivalent to the
abstract candidates of Vol~IV's
\Cref{v4-arith:rh:hyp:HPAr}, not independently constructed. The
transfer does not shrink the problem; it renames it.
\end{remark}

\section{Meyer 2005 Partial Approach and Its Limitation}
\label{v4-arith:w9-r4:meyer-2005}

Meyer's Schwartz-class refinement provides the most explicit
existing operator-level treatment. It closes the spectral-flow
problem on a dense Schwartz subspace, and transfers the residual to
an explicitly named spectral-matching conjecture.

\subsection{The Meyer Construction}
\label{v4-arith:w9-r4:meyer-2005:construction}

\begin{definition}[Meyer Schwartz space and scaling]
\label{v4-arith:w9-r4:def:meyer-data}
\ClaimStatusDefinitional \emph{(Meyer 2005 arXiv:math/0311468 \S3)}.
The \emph{Meyer Schwartz space} $\mathcal{S}_{\mathrm{Meyer}}(\mathbb{A})$
is the sub-Schwartz space of
$L^{2}(\mathbb{A}^{\times}/\mathbb{Q}^{\times},d^{\times}x)$
consisting of functions with rapid decay at all places adelically
(Meyer \S3.1). The Meyer scaling $D_{\mathrm{Meyer}}$ is the
restriction of the idele scaling $D_{0}$ of
\Cref{v4-arith:w7-j:def:H-idele} to $\mathcal{S}_{\mathrm{Meyer}}(\mathbb{A})$.
\end{definition}

\begin{theorem}[Meyer spectral-matching intertwiner]
\label{v4-arith:w9-r4:thm:meyer-intertwiner}
\ClaimStatusProvedElsewhere \emph{(Meyer 2005 Thm.~4.4, \S4.2)}. There
is an explicit Mellin-unitary intertwiner
\[
U_{\mathrm{Meyer}}:\;\mathcal{S}_{\mathrm{Meyer}}(\mathbb{A})\;\to\;L^{2}(\mathbb{R}_{1/2+it},d\tau)
\]
conjugating $D_{\mathrm{Meyer}}$ to the multiplication-by-$\tau$
operator on the critical line and landing on the Paley-Wiener
Schwartz-like class with distributional pole data at the non-trivial
zeros of $\xi_{\mathbb{R}}$.
\end{theorem}

\subsection{The Limitation: Pole Data Matching}
\label{v4-arith:w9-r4:meyer-2005:limitation}

\begin{remark}[Meyer's intertwiner is not $A$]
\label{v4-arith:w9-r4:rem:meyer-not-A}
$U_{\mathrm{Meyer}}$ is a unitary from Meyer's Schwartz space into
$L^{2}$ of the critical line. It does not intertwine $N_{\mathrm{BC}}$
with $H_{\mathrm{HP}}^{\mathrm{formal}}$; it intertwines
$D_{\mathrm{Meyer}}$ with multiplication by the spectral variable.
The pole data at non-trivial zeros is distributional, not an
eigenvalue realisation. To upgrade $U_{\mathrm{Meyer}}$ to a
spectral-flow $A$, one must pass from the distributional pole
structure on $L^{2}(\mathbb{R}_{1/2+it})$ to an $\ell^{2}$-sum over
zeros, which requires the zeros to be discrete points on the critical
line: this is exactly RH.
\end{remark}

\begin{theorem}[Meyer equivalence to H-P$^{\mathrm{Ar}}$]
\label{v4-arith:w9-r4:thm:meyer-eq-HPAr}
\ClaimStatusProvedElsewhere \emph{(Meyer 2005 \S5 conditional;
restated here)}. The Meyer spectral-matching conjecture (that
$U_{\mathrm{Meyer}}$ extends continuously to a unitary intertwiner
delivering the zeta zeros as absorption spectrum) is equivalent to
the Connes absorption-spectrum conjecture
(\Cref{v4-arith:w9-r4:thm:connes-absorption}), hence to
H-P$^{\mathrm{Ar}}$.
\end{theorem}

\begin{remark}[Meyer gives Schwartz-class closure, no more]
\label{v4-arith:w9-r4:rem:meyer-schwartz-only}
The Meyer construction closes the spectral flow unconditionally on
the Schwartz sub-space $\mathcal{S}_{\mathrm{Meyer}}(\mathbb{A})$,
intertwining $D_{\mathrm{Meyer}}$ with multiplication by $\tau$ on
$L^{2}(\mathbb{R}_{1/2+it})$. It does \emph{not} close the
identification of the pole data on the critical line with the
$\ell^{2}$-sum over zeros; that extension is the Meyer spectral-matching
conjecture, equivalent to Connes absorption, equivalent to
H-P$^{\mathrm{Ar}}$. Meyer therefore supplies the left half of the
spectral flow (Mellin transform) unconditionally and leaves the
right half (Hadamard inversion as spectral realisation) as the
conjectural kernel.
\end{remark}

\section{Residual: Irreducible Operator-Theoretic Obstruction}
\label{v4-arith:w9-r4:residual}

We now isolate the irreducible residual and name it precisely.

\subsection{Statement of the Irreducible Obstruction}
\label{v4-arith:w9-r4:residual:statement}

\begin{theorem}[irreducible residual for spectral flow]
\label{v4-arith:w9-r4:thm:irreducible-residual}
\ClaimStatusProvedHere. The spectral-flow problem
(\Cref{v4-arith:w9-r4:def:spectral-flow}) reduces to the following
operator-theoretic statement on the idele class Hilbert space
$\mathcal{H}^{\mathrm{id}}=L^{2}(C_{\mathbb{Q}},d^{\times}x)$.
\begin{quote}
\emph{A-TP$^{\mathrm{arch}}$}: The Meyer scaling generator
$D_{\mathrm{Meyer}}$, extended from the Schwartz sub-space
$\mathcal{S}_{\mathrm{Meyer}}(\mathbb{A})$ to the cuspidal subspace
$\mathcal{E}\subset\mathcal{H}^{\mathrm{id}}$, has pure point
absorption spectrum equal to
$\{\gamma_{\rho}\mid\rho\in\mathcal{Z}\}$ with multiplicities
matching. Equivalently, the closure of $U_{\mathrm{Meyer}}$ lifts
the distributional Paley--Wiener class on
$L^{2}(\mathbb{R}_{1/2+it})$ to an $\ell^{2}$-direct sum over the
non-trivial zeros.
\end{quote}
The following reductions are unconditional.
\begin{enumerate}
\item The source-side data
$(\mathcal{H}_{\mathrm{BC}},N_{\mathrm{BC}})$ is unconditional
(\Cref{v4-arith:w5-b:thm:N-BC-self-adjoint}).
\item The Meyer Mellin unitary
$U_{\mathrm{Meyer}}:\mathcal{S}_{\mathrm{Meyer}}(\mathbb{A})\to L^{2}(\mathbb{R}_{1/2+it})$
is unconditional
(\Cref{v4-arith:w9-r4:thm:meyer-intertwiner}).
\item The formal Mellin--Hadamard residue calculation is
unconditional and obstructs the naive symbolic intertwining:
Mellin scaling differentiates the numerator, while zero-ordinate
multiplication would multiply the residue
(\Cref{v4-arith:w9-r4:lem:formal-intertwine}).
\item The equivalence
H-P$^{\mathrm{Ar}}$ $\Leftrightarrow$ Connes absorption $\Leftrightarrow$
Meyer spectral-matching $\Leftrightarrow$ A-TP$^{\mathrm{arch}}$ is
proven within the Vol~IV wave-6 / wave-7 ledger
(\Cref{v4-arith:hpvbconv:thm:HPAr-equivalence},
\Cref{v4-arith:w7-j:thm:ATP-iff-S-unitary}).
\end{enumerate}
What is \emph{not} unconditional is A-TP$^{\mathrm{arch}}$ itself:
it is the absorption-spectrum property, not derivable from the
source-side data or the Mellin unitary by any known argument.
\end{theorem}

\begin{proof}
Items (1)--(4) are cited. The content of the theorem is the
reduction statement: the spectral-flow problem in
\Cref{v4-arith:w9-r4:def:spectral-flow} admits a decomposition
\[
A\;=\;(\text{absorption-spectrum inclusion})\circ U_{\mathrm{Meyer}}\circ(\text{embedding }\mathcal{H}_{\mathrm{BC}}\to\mathcal{S}_{\mathrm{Meyer}}(\mathbb{A}))
\]
in which the first two factors are unconditional and the third
factor (the absorption-spectrum inclusion into an $\ell^{2}$-direct
sum) is the irreducible residual. The embedding
$\mathcal{H}_{\mathrm{BC}}\hookrightarrow\mathcal{S}_{\mathrm{Meyer}}(\mathbb{A})$
is the adelic version of the Bost--Connes to idele scaling
embedding (Meyer 2005 \S4.1); it is bounded after the
Meyer normalisation. Under the embedding, $N_{\mathrm{BC}}$ maps to
a summand of $D_{\mathrm{Meyer}}$ restricted to the spherical
character sector; the summand corresponds to the finite-prime part
of the scaling, while the archimedean part is the Arakelov scaling
$\partial_{\log t_{\infty}}$ at the real place. The intertwining
\eqref{v4-arith:w9-r4:eq:intertwine} holds on the embedded core by
\Cref{v4-arith:w9-r4:thm:meyer-intertwiner}. The spectral identity
\eqref{v4-arith:w9-r4:eq:spec-id} requires the absorption-spectrum
property (A-TP$^{\mathrm{arch}}$), which is the Connes-Meyer
spectral-matching conjecture.
\end{proof}

\begin{remark}[the residual is at the archimedean place]
\label{v4-arith:w9-r4:rem:residual-at-infty}
The Meyer intertwining
(\Cref{v4-arith:w9-r4:thm:meyer-intertwiner}) is derived
place-by-place from the Tate 1950 local functional equations. At
each finite prime, the local spectral matching is elementary: the
BC generator $\mu_{p}$ acts with eigenvalue $p^{-it}$ on the Mellin
side, conjugate to multiplication by $t$ by a change of variable.
At the archimedean place, the Arakelov scaling acts with eigenvalue
$e^{-it}$ on the Mellin side, and the spectral identification
requires the scaling to have pure point absorption spectrum
$\{\gamma_{\rho}\}$ on the cuspidal sub-space. The latter is
equivalent to A-TP
(\Cref{v4-arith:w5-a:thm:ATP-equals-weil-positivity}) by
\Cref{v4-arith:w7-j:thm:ATP-iff-S-unitary}, hence the label
A-TP$^{\mathrm{arch}}$. The obstruction is purely archimedean; the
finite-prime spectral matching is delivered by Bost--Connes
unconditionally.
\end{remark}

\subsection{Spectral-Flow Status Ledger}
\label{v4-arith:w9-r4:residual:ledger}

\begin{center}
\small
\begin{tabular}{p{0.42\textwidth}|p{0.25\textwidth}|p{0.28\textwidth}}
Component & Status & Reference \\
\hline
Source $(\mathcal{H}_{\mathrm{BC}},N_{\mathrm{BC}})$, self-adjoint & Unconditional & \Cref{v4-arith:w5-b:thm:N-BC-self-adjoint} \\
Source Mellin inversion on $\mathcal{H}_{\mathrm{BC}}$ & Obstructed (atomic measures) & \Cref{v4-arith:w9-r4:lem:no-mellin-on-HBC} \\
Embedding $\mathcal{H}_{\mathrm{BC}}\hookrightarrow\mathcal{S}_{\mathrm{Meyer}}(\mathbb{A})$ & Unconditional after Meyer renormalisation & \Cref{v4-arith:w7-j:rem:trace-class-scope} \\
Meyer unitary $U_{\mathrm{Meyer}}$ on Schwartz sub-space & Unconditional & \Cref{v4-arith:w9-r4:thm:meyer-intertwiner} \\
Hadamard kernel $K_{\mathrm{Had}}$ Hilbert-Schmidt & Obstructed (divergent at zeros) & \Cref{v4-arith:w9-r4:lem:had-not-HS} \\
Naive formal symbolic intertwining & Obstructed by derivative residue & \Cref{v4-arith:w9-r4:lem:formal-intertwine} \\
Absorption spectrum of $D_{\mathrm{Meyer}}|_{\mathcal{E}}$ = $\{\gamma_{\rho}\}$ & Conjectural (A-TP$^{\mathrm{arch}}$) & \Cref{v4-arith:w9-r4:thm:irreducible-residual} \\
Connes equivalence with H-P$^{\mathrm{Ar}}$ & Proved equivalent & \Cref{v4-arith:w9-r4:prop:connes-eq-HPAr} \\
Meyer spectral-matching equivalence with A-TP$^{\mathrm{arch}}$ & Proved equivalent & \Cref{v4-arith:w9-r4:thm:meyer-eq-HPAr} \\
Existence of $A$ as closed operator & Conjectural = A-TP$^{\mathrm{arch}}$ & \Cref{v4-arith:w9-r4:thm:irreducible-residual} \\
Existence of $A$ as bounded operator & Impossible (see unboundedness) & \Cref{v4-arith:w9-r4:rem:unboundedness-intrinsic}
\end{tabular}
\end{center}

\subsection{Why the Residual is Irreducible}
\label{v4-arith:w9-r4:residual:why-irreducible}

\begin{theorem}[irreducibility of A-TP$^{\mathrm{arch}}$]
\label{v4-arith:w9-r4:thm:irreducibility}
\ClaimStatusProvedHere. A-TP$^{\mathrm{arch}}$ cannot be reduced to
any further spectral-flow identity on the Bost--Connes / idele
framework: it is the terminus of the reduction chain.
\end{theorem}

\begin{proof}
Under any Hilbert-space completion $\mathcal{H}_{\xi}$ of the target,
the spectrum of $D_{\mathrm{Meyer}}|_{\mathcal{E}}$ is determined
by the Tate-Mellin transform of the cuspidal forms modulo the
residue structure at non-trivial zeros. The pole-to-point spectrum
passage is precisely the spectral-matching conjecture; any finer
reduction would require an \emph{a priori} realisation of the
cuspidal forms as an $\ell^{2}$-sum over non-trivial zeros, which
presupposes the critical-line placement. No unconditional route
exists: every known partial construction (Connes' adele class,
Meyer's Schwartz refinement, de Branges' Hermite-Biehler,
Nyman-Beurling, Guinand-Weil, Bost--Connes GNS, scattering
reformulations of Chapter~\ref{v4-arith:w7-j:ATP-adelic-scattering})
transfers the problem to an equivalent absorption or positivity
statement, each equivalent to RH.

A direct obstruction: the naive symbolic intertwining
$A^{\mathrm{sym}}\,N_{\mathrm{BC}}=H_{\mathrm{HP}}\,A^{\mathrm{sym}}$
does not follow from residue calculus. By
\Cref{v4-arith:w9-r4:lem:formal-intertwine}, the left hand side
extracts the residue of \(-F'/\xi_{\mathbb R}\), while the right hand
side multiplies the residue of \(F/\xi_{\mathbb R}\) by
\(\gamma_{\rho}\). Equality requires the first-order interpolation
condition \(F'(\rho)=-\gamma_{\rho}F(\rho)\) at every non-trivial
zero. RH alone does not give this. Thus any closed-operator upgrade
of $A^{\mathrm{sym}}$ must add a genuine domain/boundary theorem.
The irreducible residual is this passage; it has no known proof
independent of RH.
\end{proof}

\begin{remark}[Beilinson-principled closure of this chapter]
\label{v4-arith:w9-r4:rem:beilinson-closure}
The honest content of the chapter is:
\begin{enumerate}
\item The source $(\mathcal{H}_{\mathrm{BC}},N_{\mathrm{BC}})$ is
unconditional.
\item The Meyer Mellin unitary on Schwartz sub-spaces is
unconditional.
\item The formal spectral flow $A^{\mathrm{sym}}$ is a well-defined
symbolic composition.
\item The symbolic-to-operatorial upgrade is the Connes-Meyer-de
Branges-Nyman irreducible obstruction, which is equivalent to RH.
\end{enumerate}
No closure of H-P$^{\mathrm{Ar}}$ is claimed. The chapter exists as
a location document, naming the irreducible obstruction precisely
within the Vol~IV framework.
\end{remark}

\section{Placement in Triple-Identification Framework}
\label{v4-arith:w9-r4:placement}

We close by placing the spectral-flow problem in the triple
identification of Chapters~\ref{v4-arith:w8-a:deninger-chiral-homology},
\ref{v4-arith:w7-j:ATP-adelic-scattering}, and
\ref{v4-arith:3d-acs-E1bar}.

\subsection{Spectral Flow as Deninger Spectral-Flow Theorem}
\label{v4-arith:w9-r4:placement:deninger}

\begin{proposition}[spectral flow = Deninger spectral-flow]
\label{v4-arith:w9-r4:prop:spectral-flow-eq-deninger}
\ClaimStatusConditional \emph{(on the triple-identification and
\(L_0\)-conjugacy hypotheses).} Under the Deninger triple
identification and the \(L_0\)-conjugacy hypotheses cited below, the
spectral-flow problem of \Cref{v4-arith:w9-r4:def:spectral-flow}
coincides with the Deninger spectral-flow conjecture
\Cref{v4-arith:3d-acs:conj:spectral-flow}. This is a dictionary
identification of two names for the same residual, not an independent
construction of the spectral-flow operator.
\end{proposition}

\begin{proof}
Unpack the Deninger conjecture
\Cref{v4-arith:3d-acs:conj:spectral-flow}: it posits a closed
operator $\Theta_{\mathrm{Den}}=\tfrac12\mathrm{id}+iH_{\mathrm{Den}}$
on a completion of the primary middle cohomology
$H^{1}_{\mathrm{Den},\chi,\mathrm{prim}}$ with $H_{\mathrm{Den}}$
self-adjoint in the Petersson/GNS inner product, determinant
$\xi_{\mathbb{R}}$, and Weil--Connes trace distribution equal to the
zero-side. The Hilbert completion is
$\mathcal{H}_{\xi}^{\mathrm{Petersson}}$ of
\Cref{v4-arith:w9-r4:def:target-completion}, and the operator
$H_{\mathrm{Den}}$ is $H_{\mathrm{HP}}$ in the Vol~IV vocabulary
(\Cref{v4-arith:rh:rem:spectral-convention}).

The Deninger triple identification
(\Cref{v4-arith:w8-a:thm:triple-id}) identifies the Petersson
completion with the arithmetic chiral homology of
$\mathsf{H}^{\mathrm{Ar}}$ and with the boundary chiral homology of
$\mathrm{ACS}^{\mathrm{qu}}_{k^{\mathrm{Ar}}}$. The spectral-flow
operator $A$ of \Cref{v4-arith:w9-r4:def:spectral-flow} is exactly
the unitary $U$ of
\Cref{v4-arith:w8-a:lem:L0-conjugacy}: it conjugates the archimedean
Frobenius $F_{\infty}=e^{-\partial_{t}}$ to the Virasoro-$L_{0}$
representation on the arithmetic Heisenberg Fock, and via the primary-Fock
embedding of \Cref{v4-arith:w5-b:lem:fock-embeds-gns} maps
$N_{\mathrm{BC}}$ to the zero-side operator $H_{\mathrm{HP}}$ on
the completed primary sector.
\end{proof}

\begin{remark}[the Deninger framework renames the residual]
\label{v4-arith:w9-r4:rem:deninger-renames}
The Deninger framework does not close the spectral-flow problem;
it labels the same residual with the arithmetic-cohomology name.
The irreducible obstruction A-TP$^{\mathrm{arch}}$ of
\Cref{v4-arith:w9-r4:thm:irreducible-residual} is Deninger axiom
(3) (regularised-determinant identity) in the Vol~IV realisation
(\Cref{v4-arith:w8-a:thm:deninger-chiral-equals}, conditional on
the spectral-flow conjecture). Naming the problem in cohomological
language does not shrink it.
\end{remark}

\subsection{Spectral Flow as A-TP Adelic-Scattering}
\label{v4-arith:w9-r4:placement:atp}

\begin{proposition}[spectral flow = A-TP adelic-scattering]
\label{v4-arith:w9-r4:prop:spectral-flow-eq-atp}
\ClaimStatusProvedHere. The irreducible residual
A-TP$^{\mathrm{arch}}$ of
\Cref{v4-arith:w9-r4:thm:irreducible-residual} coincides with A-TP
on $\mathrm{PW}(\mathbb{R})\setminus\mathrm{PW}^{\mathrm{TM}}$ in
the scattering formulation of
Chapter~\ref{v4-arith:w7-j:ATP-adelic-scattering}.
\end{proposition}

\begin{proof}
The Meyer spectral-matching conjecture
(\Cref{v4-arith:w9-r4:thm:meyer-eq-HPAr}) is equivalent to unitarity
of the scattering matrix $S$ of the Bost--Connes-perturbed scaling
generator pair $(D_{0},D_{\mathrm{BC}})$ on the full
$\mathrm{PW}(\mathbb{R})$; this is exactly
\Cref{v4-arith:w7-j:thm:ATP-iff-S-unitary} in its forward direction.
The extension from Tate-Mellin sub-class
$\mathrm{PW}^{\mathrm{TM}}$ to the full $\mathrm{PW}(\mathbb{R})$ is
\Cref{v4-arith:w7-j:thm:obstruction-located}. Under the embedding
$\mathcal{H}_{\mathrm{BC}}\hookrightarrow\mathcal{S}_{\mathrm{Meyer}}(\mathbb{A})$
of the spectral-flow setup, the absorption-spectrum condition
becomes the scattering-matrix unitarity condition. Identification.
\end{proof}

\begin{remark}[scattering transfers the obstruction]
\label{v4-arith:w9-r4:rem:scattering-transfers}
The scattering formulation of
Chapter~\ref{v4-arith:w7-j:ATP-adelic-scattering} closes the
spectral flow unconditionally on the Tate-Mellin sub-class
$\mathrm{PW}^{\mathrm{TM}}$ and transfers the residual on
$\mathrm{PW}(\mathbb{R})\setminus\mathrm{PW}^{\mathrm{TM}}$ to the
Connes non-commutative $L^{2}$-structure on the adele class space.
In spectral-flow vocabulary, Tate-Mellin closure corresponds to
the zero-adapted Schwartz functions on $\mathcal{H}_{\mathrm{BC}}$
where the Hadamard inversion is algebraically realisable by local
functional equations; the full extension requires the absorption
property.
\end{remark}

\subsection{Spectral Flow as 3D ACS Bulk-Boundary Correspondence}
\label{v4-arith:w9-r4:placement:acs}

\begin{proposition}[spectral flow = 3D ACS level-deformation]
\label{v4-arith:w9-r4:prop:spectral-flow-eq-acs}
\ClaimStatusProvedHereConditional \emph{(on
\Cref{v4-arith:3d-acs:conj:qu-acs})}. Under the 3D arithmetic
Chern-Simons conjecture
\Cref{v4-arith:3d-acs:conj:qu-acs}, the spectral-flow operator $A$ is
the level-derivative operator on the 3D bulk Hilbert space
$\mathcal{V}^{\mathrm{prim}}$.
\end{proposition}

\begin{proof}
The 3D ACS bulk partition function satisfies
$Z_{\mathrm{ACS}^{\mathrm{qu}}}(\overline{\mathrm{Spec}\,\mathbb{Z}};s)=\xi_{\mathbb{R}}(s)$
(\Cref{v4-arith:3d-acs:cor:partition-bar}). The level-derivative
operator $\partial_{s}$ on the bulk Hilbert space acts, by the
canonical-quantisation axiom
\Cref{v4-arith:3d-acs:conj:qu-acs}(2), as the derivative of the
Virasoro-$L_{0}$ eigenvalue with respect to the spectral parameter,
equivalently as the Arakelov-scaling Frobenius $F_{\infty}$ of
\Cref{v4-arith:w8-a:def:F-infinity}. The spectral-flow operator $A$
conjugates this to the multiplication-by-$\gamma_{\rho}$ operator on
the zero index set via the Mellin unitary of
\Cref{v4-arith:w8-a:lem:L0-conjugacy}.
\end{proof}

\begin{remark}[the 3D framework renames the obstruction as Wilson-loop matching]
\label{v4-arith:w9-r4:rem:acs-wilson}
The 3D ACS framework reformulates the spectral-flow obstruction as
a Wilson-loop correlator matching condition: the level-derivative
operator's spectrum must match the zero multiset under the
Arakelov-K\"ahler Wilson-surface realisation
(\Cref{v4-arith:w8-a:rem:wilson-infinity}). This is equivalent to
Deninger axiom (2), to Meyer spectral-matching, to Connes absorption,
to A-TP$^{\mathrm{arch}}$, to H-P$^{\mathrm{Ar}}$. The 3D framework
adds geometric content (the bulk theory becomes a concrete
Chern-Simons computation) but does not change the obstruction's
strength.
\end{remark}

\subsection{The Unified Residual}
\label{v4-arith:w9-r4:placement:unified}

\begin{theorem}[unified residual, five equivalent names]
\label{v4-arith:w9-r4:thm:unified-residual}
\ClaimStatusProvedHere \emph{(conditional on the referenced
equivalences)}. The spectral-flow problem
(\Cref{v4-arith:w9-r4:def:spectral-flow}) reduces to a single
operator-theoretic statement, carrying five equivalent names across
the Vol~IV framework:
\begin{enumerate}
\item \emph{H-P$^{\mathrm{Ar}}$} (\Cref{v4-arith:rh:hyp:HPAr}):
existence of a self-adjoint operator $H_{\mathrm{HP}}$ with
spectrum $\{\gamma_{\rho}\}$ on a Petersson completion.
\item \emph{Connes absorption}
(\Cref{v4-arith:w9-r4:thm:connes-absorption}): absorption spectrum
of the scaling generator on the cuspidal subspace of the adele
class space.
\item \emph{Meyer spectral-matching}
(\Cref{v4-arith:w9-r4:thm:meyer-eq-HPAr}): extension of the Meyer
Mellin unitary from Schwartz sub-spaces to the full cuspidal sector.
\item \emph{A-TP} (\Cref{v4-arith:w5-a:conj:archimedean-tate-positivity}):
Weil-positivity of the archimedean Tate distribution on
$\mathrm{PW}(\mathbb{R})$.
\item \emph{Deninger spectral-flow}
(\Cref{v4-arith:3d-acs:conj:spectral-flow}): existence of the
middle-cohomology operator in the arithmetic cohomology of
$\overline{\mathrm{Spec}\,\mathbb{Z}}$.
\end{enumerate}
All five are equivalent to each other and to the Riemann Hypothesis
for $\zeta$.
\end{theorem}

\begin{proof}
Within the Vol~IV framework, the five equivalences are recorded as:
(1) $\Leftrightarrow$ (2) by
\Cref{v4-arith:w9-r4:prop:connes-eq-HPAr};
(2) $\Leftrightarrow$ (3) by
\Cref{v4-arith:w9-r4:thm:meyer-eq-HPAr};
(3) $\Leftrightarrow$ (4) by
\Cref{v4-arith:w7-j:thm:ATP-iff-S-unitary};
(4) $\Leftrightarrow$ RH by
\Cref{v4-arith:w5-a:thm:ATP-equals-weil-positivity};
(1) $\Leftrightarrow$ (5) by
\Cref{v4-arith:w9-r4:prop:spectral-flow-eq-deninger}. The chain
closes with all five mutually equivalent to RH.
\end{proof}

\begin{remark}[honest summary]
\label{v4-arith:w9-r4:rem:honest-summary}
The spectral-flow problem $N_{\mathrm{BC}}\rightsquigarrow H_{\mathrm{HP}}$
is the classical Hilbert--P\'olya twist. It has been recast
five ways within the Vol~IV framework; each recast supplies
different geometric or analytic content and closes the problem on
different explicit sub-classes, but none closes the residual. The
residual is irreducible within the Vol~IV setting: it is the
absorption-spectrum property of the archimedean scaling on the
cuspidal subspace, equivalent to RH.

The Vol~IV contribution is \emph{location}. The chapter catalogues
the unconditional source, the unconditional Mellin unitary on
Schwartz sub-spaces, the formal Mellin--Hadamard obstruction, and names
the irreducible operator-theoretic obstruction precisely. It does
not close H-P$^{\mathrm{Ar}}$. One hundred years after Hilbert and
P\'olya, the operator construction remains open; the present
chapter exhibits the Vol~IV packaging of the open residual without
pretending closure.
\end{remark}

\section{Polyakov--Dirac Normal Form for the Residual}
\label{v4-arith:w9-r4:polyakov-dirac-normal-form}

The two useful reductions of the previous sections are orthogonal.
The Polyakov reduction fixes the anomaly coefficient; the Dirac
reduction forces the remaining problem into a first-order
self-adjointness statement. Together they remove the two common
false escape routes: moving the central charge, and hiding the
Hilbert--P\'olya content inside an unspecified ``flow''.

\subsection{Polyakov Rigidity of the Scalar}
\label{v4-arith:w9-r4:pdnf:polyakov}

\begin{proposition}[Polyakov scalar rigidity]
\label{v4-arith:w9-r4:prop:polyakov-scalar-rigidity}
\ClaimStatusProvedHere. In the primary Heisenberg scope of
Chapters~\ref{v4-arith:polyakov} and
\ref{v4-arith:w9-r3:chapter}, any spectral-flow closure compatible
with
\begin{enumerate}
\item the rank-$2$ per-prime Sugawara Heisenberg stress tensor,
\item the global Fock partition
$Z(s)=\prod_{p}\prod_{n\geq 1}(1-p^{-sn})^{-2}$, and
\item the archimedean Polyakov anomaly of the rank-$2$ free-field
completion,
\end{enumerate}
has scalar Virasoro anomaly
\begin{equation}
\label{v4-arith:w9-r4:eq:polyakov-rigid-c}
c_{\mathrm{reg}}\;=\;2.
\end{equation}
Consequently the residual spectral-flow problem cannot be shifted
into a different central-charge normalisation. Its remaining
content is the operator-domain and positivity problem named
A-TP$^{\mathrm{arch}}$ in
\Cref{v4-arith:w9-r4:thm:irreducible-residual}.
\end{proposition}

\begin{proof}
\Cref{v4-arith:w9-r3:thm:main} proves that the naive scalar
$\sum_{p}c^{(p)}$ diverges, that the per-state Fock action is
finite only after separating off the universal scalar term, and
that the global Fock-partition residue gives
$c_{\mathrm{reg}}=2$. The same chapter proves the Polyakov
regularisation identity
\Cref{v4-arith:w9-r3:thm:polyakov-identity}, identifying this
finite-prime residue with the rank-$2$ archimedean conformal
anomaly. These inputs leave no scalar parameter: changing
$c_{\mathrm{reg}}$ would either abandon the Euler-product
partition, abandon the Sugawara rank, or abandon the archimedean
anomaly match. None of those changes is a spectral-flow closure
inside the stated primary Heisenberg scope. The only remaining
unsettled part is therefore the closed-operator passage from the
Meyer/Connes absorption symbol to a self-adjoint zero operator,
which is exactly A-TP$^{\mathrm{arch}}$ by
\Cref{v4-arith:w9-r4:thm:irreducible-residual}.
\end{proof}

\begin{remark}[what Polyakov does and does not buy]
\label{v4-arith:w9-r4:rem:polyakov-no-rh}
Polyakov regularisation is decisive for the anomaly bookkeeping,
not for RH. It supplies the finite scalar $c_{\mathrm{reg}}=2$ and
rules out the divergent prime-zeta prescription; it does not
produce the Hilbert space on which the zero ordinates are a
self-adjoint spectrum.
\end{remark}

\subsection{Dirac Closure of the Remaining Operator}
\label{v4-arith:w9-r4:pdnf:dirac}

\begin{definition}[residual symmetric ordinate core]
\label{v4-arith:w9-r4:def:residual-ordinate-core}
\ClaimStatusDefinitional. A \emph{residual symmetric ordinate core}
is data
\[
(\mathcal{C}_{\mathrm{res}},H_{\mathrm{res}}^{\mathrm{sym}})
\]
where \(\mathcal{C}_{\mathrm{res}}\subset\mathcal{E}\) is an
\(L^{2}\)-dense candidate core and \(H_{\mathrm{res}}^{\mathrm{sym}}\)
is a symmetric operator on it. The intended source of
\(\mathcal{C}_{\mathrm{res}}\) is the Meyer Mellin unitary applied to
the Schwartz adelic core, together with a Hadamard-residue boundary
condition at the non-trivial zeros, as in
\Cref{v4-arith:w9-r4:thm:meyer-intertwiner} and
\Cref{v4-arith:w9-r4:def:formal-MH}. This sentence is a source
description, not a graph-density theorem. The graph-density and
non-tautology of the Hadamard-residue boundary condition are the
residual obligations isolated below.

On the Mellin line \(H_{\mathrm{res}}^{\mathrm{sym}}\) is represented
by multiplication by the real spectral variable \(t\). After the
Hadamard-residue boundary condition it is intended to become formal
multiplication by \(\gamma_{\rho}\) on the zero components. Its
closure is not assumed to exist as a self-adjoint operator. The
superscript ``\(\mathrm{sym}\)'' is part of the obligation: if the
residue data produces non-real ordinate coefficients, no symmetric
operator has yet been constructed, and the Dirac closure below fails.
\end{definition}

\begin{definition}[Dirac double of the residual]
\label{v4-arith:w9-r4:def:dirac-double}
\ClaimStatusDefinitional. The \emph{Dirac double} of the residual
ordinate core is the odd symmetric operator
\begin{equation}
\label{v4-arith:w9-r4:eq:dirac-double}
\mathscr{D}_{\mathrm{res}}^{0}
\;:=\;
\begin{pmatrix}
0 & H_{\mathrm{res}}^{\mathrm{sym}}\\
H_{\mathrm{res}}^{\mathrm{sym}} & 0
\end{pmatrix}
\quad\text{on}\quad
\mathcal{C}_{\mathrm{res}}\oplus\mathcal{C}_{\mathrm{res}},
\end{equation}
with grading
$\Gamma=\operatorname{diag}(1,-1)$. A \emph{Dirac closure} is a
self-adjoint closure
$\mathscr{D}_{\mathrm{res}}=\overline{\mathscr{D}_{\mathrm{res}}^{0}}$
whose spectral measure satisfies the regularised determinant
identity
\begin{equation}
\label{v4-arith:w9-r4:eq:dirac-det}
\det\nolimits_{\infty}\bigl(s\cdot\mathrm{id}
-({\tfrac12}\mathrm{id}+iH_{\mathrm{res}})\bigr)
\;=\;\xi_{\mathbb{R}}(s),
\end{equation}
where $H_{\mathrm{res}}$ is the off-diagonal component recovered
from $\mathscr{D}_{\mathrm{res}}$ by the grading.
\end{definition}

\begin{lemma}[Dirac double detects self-adjointness]
\label{v4-arith:w9-r4:lem:dirac-double-detects-sa}
\ClaimStatusProvedHere. Let $H$ be a densely defined symmetric
operator on a Hilbert space $\mathcal{H}$ and
$\mathscr{D}_{H}^{0}=
\left(\begin{smallmatrix}0&H\\H&0\end{smallmatrix}\right)$
on $\mathcal{D}(H)\oplus\mathcal{D}(H)$. Then
$\mathscr{D}_{H}^{0}$ is essentially self-adjoint if and only if
$H$ is essentially self-adjoint. When this holds,
\[
\overline{\mathscr{D}_{H}^{0}}
\;=\;
\begin{pmatrix}
0 & \overline{H}\\
\overline{H} & 0
\end{pmatrix},
\qquad
\overline{\mathscr{D}_{H}^{0}}^{\,2}
\;=\;
\begin{pmatrix}
\overline{H}^{\,2} & 0\\
0 & \overline{H}^{\,2}
\end{pmatrix}.
\]
\end{lemma}

\begin{proof}
For a symmetric $H$,
$(\mathscr{D}_{H}^{0})^{*}=
\left(\begin{smallmatrix}0&H^{*}\\H^{*}&0\end{smallmatrix}\right)$ on
$\mathcal{D}(H^{*})\oplus\mathcal{D}(H^{*})$. Hence the closure of
$\mathscr{D}_{H}^{0}$ is self-adjoint exactly when the closure of
$H$ has adjoint domain equal to its own domain, i.e.\ exactly when
$H$ is essentially self-adjoint. The displayed square is the usual
block-matrix product on the common graph domain, justified after
closure by the spectral theorem for the self-adjoint operator
$\overline{H}$ (Reed--Simon I, Ch.~VIII).
\end{proof}

\begin{theorem}[Polyakov--Dirac normal form]
\label{v4-arith:w9-r4:thm:polyakov-dirac-normal-form}
\ClaimStatusProvedHere. Within the spectral-flow framework of
\Cref{v4-arith:w9-r4:def:spectral-flow}, after imposing the
Polyakov scalar rigidity
\Cref{v4-arith:w9-r4:prop:polyakov-scalar-rigidity}, the following
statements are equivalent:
\begin{enumerate}
\item H-P$^{\mathrm{Ar}}$ for $\xi_{\mathbb{R}}$:
there exists a self-adjoint zero ordinate operator
$H_{\mathrm{HP}}$ whose regularised determinant is
$\xi_{\mathbb{R}}$.
\item A-TP$^{\mathrm{arch}}$:
the Meyer/Connes absorption spectrum on the cuspidal subspace is
the non-trivial-zero ordinate multiset.
\item The residual symmetric ordinate core
$H_{\mathrm{res}}^{\mathrm{sym}}$ of
\Cref{v4-arith:w9-r4:def:residual-ordinate-core} admits a Dirac
closure in the sense of \Cref{v4-arith:w9-r4:def:dirac-double}.
\end{enumerate}
Thus the entire remaining load-bearing block is the essential
self-adjointness and determinant identity of one first-order
Dirac double on the residual Meyer core. No central-charge,
partition-function, or formal trace adjustment can change this
equivalence.
\end{theorem}

\begin{proof}
(1) $\Leftrightarrow$ (2) is
\Cref{v4-arith:w9-r4:prop:connes-eq-HPAr} together with
\Cref{v4-arith:w9-r4:thm:meyer-eq-HPAr} and the residual theorem
\Cref{v4-arith:w9-r4:thm:irreducible-residual}.

(1) $\Rightarrow$ (3): given H-P$^{\mathrm{Ar}}$, take
$H_{\mathrm{res}}=H_{\mathrm{HP}}$ on the H-P Hilbert space and
restrict to the common Meyer/Hadamard core. The Dirac double
\eqref{v4-arith:w9-r4:eq:dirac-double} is essentially
self-adjoint by \Cref{v4-arith:w9-r4:lem:dirac-double-detects-sa},
and the determinant identity
\eqref{v4-arith:w9-r4:eq:dirac-det} is exactly the
H-P$^{\mathrm{Ar}}$ determinant identity
\Cref{v4-arith:rh:eq:HP-determinant}.

(3) $\Rightarrow$ (1): a Dirac closure gives a self-adjoint
off-diagonal component $H_{\mathrm{res}}$ by the same lemma. The
regularised determinant identity
\eqref{v4-arith:w9-r4:eq:dirac-det} states precisely that
$\tfrac12\mathrm{id}+iH_{\mathrm{res}}$ has characteristic
determinant $\xi_{\mathbb{R}}$, hence its spectral multiset is the
non-trivial-zero multiset with multiplicity. This is
H-P$^{\mathrm{Ar}}$ in the convention of
\Cref{v4-arith:rh:hyp:HPAr}. The Polyakov scalar rigidity is used
only to fix the anomaly counterterm at $c_{\mathrm{reg}}=2$; it
does not enter the self-adjointness implication and therefore
cannot be traded against it.
\end{proof}

\begin{remark}[single estimate left]
\label{v4-arith:w9-r4:rem:single-estimate-left}
The Dirac formulation gives the sharp analytic target. On the
Meyer residual core one must prove a graph-norm closure statement:
\[
\|u\|^{2}_{\mathscr{D}}
\;:=\;
\|u\|^{2}+\|H_{\mathrm{res}}^{\mathrm{sym}}u\|^{2}
\]
is complete after a closed Tate-polar operator quotient has first
been installed, and the closure has determinant
\eqref{v4-arith:w9-r4:eq:dirac-det}. This is not easier than RH;
it is the same obstruction with all auxiliary parameters removed.
\end{remark}

\begin{definition}[closed Tate-polar operator quotient]
\label{v4-arith:w9-r4:def:tate-polar-operator-quotient}
\ClaimStatusConditional. A \emph{closed Tate-polar operator quotient}
for \(H_{0}=H_{\mathrm{res}}^{\mathrm{sym}}\) is additional data of a
Hilbert quotient
\[
\pi\colon \mathcal{H}_{\mathrm{res}}\longrightarrow
\mathcal{H}_{\mathrm{res}}^{\circ}
\]
with closed finite-rank kernel \(P_{\mathrm{Tate}}\), an induced
dense domain
\[
\mathcal{C}_{\mathrm{res}}^{\circ}:=\pi(\mathcal{C}_{\mathrm{res}}),
\]
and a symmetric operator \(H_{0}^{\circ}\) on
\(\mathcal{C}_{\mathrm{res}}^{\circ}\) satisfying
\[
H_{0}^{\circ}\pi u=\pi H_{0}u,\qquad u\in\mathcal{C}_{\mathrm{res}}.
\]
The quotient is required to be compatible with closure and adjoint
domains:
\[
\mathcal{D}(\overline{H_{0}^{\circ}})
=\pi\mathcal{D}(\overline{H_{0}}),
\qquad
\mathcal{D}((H_{0}^{\circ})^{*})
=\pi\mathcal{D}(H_{0}^{*}),
\]
after completing both sides in their graph norms. It is also required
to be compatible with Meyer coordinates, i.e. \(U_{\mathrm{Meyer}}\)
descends to a unitary residual chart on the quotient.

The present chapter does not construct this quotient. Every phrase
``after quotienting the Tate polar directions'' in the residual
operator argument is to be read under this hypothesis.
\end{definition}

\subsection{What Remains: Boundary Form and Deficiency Kernels}
\label{v4-arith:w9-r4:pdnf:what-remains}

The normal form above is useful only if it tells us what to try next.
After the Polyakov scalar has been fixed, there is no remaining
renormalisation parameter to tune. The next attack must be on the
operator boundary itself.

\begin{definition}[residual boundary form]
\label{v4-arith:w9-r4:def:residual-boundary-form}
\ClaimStatusDefinitional. Let
\[
H_{0}:=H_{\mathrm{res}}^{\mathrm{sym}}
\quad\text{on}\quad
\mathcal{C}_{\mathrm{res}}.
\]
The adjoint-domain boundary form is the Hermitian symplectic form
\begin{equation}
\label{v4-arith:w9-r4:eq:boundary-form}
\mathcal{B}_{\mathrm{res}}(u,v)
\;:=\;
\langle H_{0}^{*}u,v\rangle-\langle u,H_{0}^{*}v\rangle,
\qquad
u,v\in\mathcal{D}(H_{0}^{*}).
\end{equation}
The \emph{residual boundary space} is the quotient
\begin{equation}
\label{v4-arith:w9-r4:eq:boundary-space}
\mathcal{Q}_{\mathrm{res}}
\;:=\;
\mathcal{D}(H_{0}^{*})/\mathcal{D}(\overline{H_{0}}),
\end{equation}
with the induced form still denoted \(\mathcal{B}_{\mathrm{res}}\).
\end{definition}

\begin{proposition}[operator boundary target after scalar rigidity]
\label{v4-arith:w9-r4:prop:polyakov-boundary-target}
\ClaimStatusProvedHere. In the primary Heisenberg scope, the only
remaining operator-theoretic route to the spectral-flow closure is a
vanishing theorem for the residual boundary quotient:
\begin{equation}
\label{v4-arith:w9-r4:eq:polyakov-boundary-target}
\mathcal{Q}_{\mathrm{res}}=0.
\end{equation}
A future Feynman/Ward proof would have to identify its BV--BFV
boundary phase space with \(\mathcal{Q}_{\mathrm{res}}\) and its
boundary pairing with \(\mathcal{B}_{\mathrm{res}}\). Cancellation of
a boundary variation merely on the original core is weaker: it proves
isotropy of the core, not \(\mathcal{Q}_{\mathrm{res}}=0\).
\end{proposition}

\begin{proof}
For a densely defined symmetric operator, the obstruction to
essential self-adjointness is precisely the boundary quotient
\(\mathcal{D}(H_{0}^{*})/\mathcal{D}(\overline{H_{0}})\), together
with its Green boundary form
\eqref{v4-arith:w9-r4:eq:boundary-form}. If the quotient vanishes,
the closure has adjoint domain equal to its own domain and is
self-adjoint. If the quotient is non-zero, there remains a genuine
choice of self-adjoint extension or a failure of self-adjointness,
and the zero-side Hilbert--Polya operator has not been constructed.

\Cref{v4-arith:w9-r4:prop:polyakov-scalar-rigidity} says that the
scalar anomaly coefficient is already fixed by the rank-\(2\)
Sugawara sector and the Fock determinant. Hence a local
counterterm cannot be used to absorb a non-zero
\(\mathcal{Q}_{\mathrm{res}}\) without leaving the primary
Heisenberg scope. What remains is therefore the operator-boundary
vanishing problem. A path-integral/Ward proof would be a proof of the
same statement only after the missing BV--BFV identification has been
constructed.
\end{proof}

\begin{definition}[deficiency spaces of the residual]
\label{v4-arith:w9-r4:def:deficiency-spaces}
\ClaimStatusDefinitional. The residual deficiency spaces are
\begin{equation}
\label{v4-arith:w9-r4:eq:deficiency-spaces}
\mathcal{N}_{+}:=\ker(H_{0}^{*}-i),
\qquad
\mathcal{N}_{-}:=\ker(H_{0}^{*}+i).
\end{equation}
\end{definition}

\begin{theorem}[Dirac two-test form of the remaining block]
\label{v4-arith:w9-r4:thm:dirac-two-test}
\ClaimStatusProvedHere. After Polyakov scalar rigidity, the
remaining load-bearing block is equivalent to the following two
statements:
\begin{enumerate}
\item \emph{Deficiency vanishing:}
\begin{equation}
\label{v4-arith:w9-r4:eq:deficiency-vanishing}
\mathcal{N}_{+}=0=\mathcal{N}_{-}.
\end{equation}
\item \emph{Determinant identity for the closure:}
\begin{equation}
\label{v4-arith:w9-r4:eq:closure-det}
\det\nolimits_{\infty}
\bigl(s\cdot\mathrm{id}
-({\tfrac12}\mathrm{id}+i\overline{H_{0}})\bigr)
\;=\;
\xi_{\mathbb{R}}(s).
\end{equation}
\end{enumerate}
Equivalently, the Dirac double
\(\mathscr{D}_{\mathrm{res}}^{0}\) has no deficiency states and its
closed off-diagonal component has determinant \(\xi_{\mathbb R}\).
\end{theorem}

\begin{proof}
By von Neumann's deficiency-index criterion, a densely defined
symmetric operator \(H_{0}\) is essentially self-adjoint if and only
if
\(\ker(H_{0}^{*}-i)=\ker(H_{0}^{*}+i)=0\). Thus
\eqref{v4-arith:w9-r4:eq:deficiency-vanishing} is exactly the
essential self-adjointness part of
\Cref{v4-arith:w9-r4:def:dirac-double}. By
\Cref{v4-arith:w9-r4:lem:dirac-double-detects-sa}, the same
condition is equivalent to essential self-adjointness of the Dirac
double \(\mathscr{D}_{\mathrm{res}}^{0}\).

Once the closure exists as a self-adjoint operator, the only
remaining datum in H-P\(^{\mathrm{Ar}}\) is the regularised
determinant. Equation
\eqref{v4-arith:w9-r4:eq:closure-det} is precisely
\eqref{v4-arith:w9-r4:eq:dirac-det} for the closure
\(\overline{H_{0}}\). Combining deficiency vanishing with the
determinant identity is therefore equivalent to the Dirac closure
condition of
\Cref{v4-arith:w9-r4:thm:polyakov-dirac-normal-form}, hence to the
remaining spectral-flow block.
\end{proof}

\begin{remark}[the smartest next calculation]
\label{v4-arith:w9-r4:rem:smartest-next-calculation}
The most concentrated attack is now concrete. Compute the two
deficiency kernels
\(\ker(H_{0}^{*}\mp i)\) on the Meyer/Hadamard residual core. If they
vanish, the operator obstruction is gone and the remaining check is
the determinant identity
\eqref{v4-arith:w9-r4:eq:closure-det}. If either kernel is non-zero,
the failure is visible as a specific boundary state in
\(\mathcal{Q}_{\mathrm{res}}\). This is the Polyakov--Dirac content:
anomaly bookkeeping fixes the scalar; first-order operator theory
exposes the only surviving obstruction.
\end{remark}

\subsection{The Deficiency Calculation}
\label{v4-arith:w9-r4:pdnf:deficiency-calculation}

We first compute the deficiency kernels abstractly. This is the
piece that makes the next execution attempt sharp: in Meyer
coordinates the kernels are range defects for multiplication by
\(t\pm i\), with the sign dictated by the convention that the
inner product is conjugate-linear in the second slot. This calculation
is unconditional only after a residual Meyer chart and the
Tate-polar operator quotient of
\Cref{v4-arith:w9-r4:def:tate-polar-operator-quotient} have been
installed. Thus the obstruction is not an Euler-product calculation
but a graph-core question for the residual Hadamard domain.

\begin{definition}[Meyer-coordinate residual chart]
\label{v4-arith:w9-r4:def:meyer-coordinate-core}
\ClaimStatusDefinitional. Let
\[
M_{t}\colon \mathcal{D}(M_{t})\subset L^{2}(\mathbb{R},dt)\to
L^{2}(\mathbb{R},dt),\qquad (M_{t}f)(t)=t f(t),
\]
be the self-adjoint multiplication operator. A \emph{Meyer-coordinate
residual chart} is a unitary map, denoted \(U_{\mathrm{Meyer}}\) by
abuse of notation,
\[
U_{\mathrm{Meyer}}\colon\mathcal{H}_{\mathrm{res}}^{\circ}
\longrightarrow L^{2}(\mathbb{R},dt),
\]
from the residual quotient Hilbert space of
\Cref{v4-arith:w9-r4:def:tate-polar-operator-quotient} such that the
Meyer-coordinate residual core is
\begin{equation}
\label{v4-arith:w9-r4:eq:meyer-coordinate-core}
\mathcal{C}_{\mathrm{res}}^{M}
\;:=\;
U_{\mathrm{Meyer}}(\mathcal{C}_{\mathrm{res}})
\;\subset\;
\mathcal{D}(M_{t}),
\end{equation}
and
\[
H_{0}=H_{\mathrm{res}}^{\mathrm{sym}}
\quad\text{is represented by}\quad
M_{t}|_{\mathcal{C}_{\mathrm{res}}^{M}} .
\]
The unconditional Meyer theorem
\Cref{v4-arith:w9-r4:thm:meyer-intertwiner} supplies this chart only
on the Schwartz adelic subspace. Its extension to the residual
quotient and to adjoint domains is part of the residual graph-core
obligation.
\end{definition}

\begin{theorem}[deficiency kernels are range defects]
\label{v4-arith:w9-r4:thm:deficiency-range-defect}
\ClaimStatusProvedHere. Assume a Meyer-coordinate residual chart in
the sense of \Cref{v4-arith:w9-r4:def:meyer-coordinate-core}. Then
\begin{equation}
\label{v4-arith:w9-r4:eq:deficiency-range-defect-plus}
U_{\mathrm{Meyer}}\mathcal{N}_{+}
\;=\;
\bigl((M_{t}+i)\mathcal{C}_{\mathrm{res}}^{M}\bigr)^{\perp},
\end{equation}
and
\begin{equation}
\label{v4-arith:w9-r4:eq:deficiency-range-defect-minus}
U_{\mathrm{Meyer}}\mathcal{N}_{-}
\;=\;
\bigl((M_{t}-i)\mathcal{C}_{\mathrm{res}}^{M}\bigr)^{\perp}.
\end{equation}
Consequently \(\mathcal{N}_{+}=0=\mathcal{N}_{-}\) if and only if
both ranges
\[
(M_{t}-i)\mathcal{C}_{\mathrm{res}}^{M},
\qquad
(M_{t}+i)\mathcal{C}_{\mathrm{res}}^{M}
\]
are dense in \(L^{2}(\mathbb{R},dt)\).
\end{theorem}

\begin{proof}
We prove the \(+\) identity; the \(-\) identity is identical with
\(i\) replaced by \(-i\). The inner product is linear in the first
slot and conjugate-linear in the second. By definition,
\(u\in\mathcal{N}_{+}\) means \(u\in\mathcal{D}(H_{0}^{*})\) and
\((H_{0}^{*}-i)u=0\). Equivalently, for every
\(f\in\mathcal{C}_{\mathrm{res}}\),
\[
\langle H_{0}f,u\rangle
\;=\;
\langle f,iu\rangle .
\]
Since the second slot is conjugate-linear,
\(\langle f,iu\rangle=-i\langle f,u\rangle\).
Transport by \(U_{\mathrm{Meyer}}\). Write
\(\tilde u=U_{\mathrm{Meyer}}u\) and
\(\tilde f=U_{\mathrm{Meyer}}f\). Since \(H_{0}\) is represented by
\(M_{t}\) on \(\mathcal{C}_{\mathrm{res}}^{M}\), the adjoint equation
becomes
\[
\langle M_{t}\tilde f,\tilde u\rangle
\;=\;
\langle \tilde f,i\tilde u\rangle ,
\qquad
\tilde f\in\mathcal{C}_{\mathrm{res}}^{M}.
\]
Hence
\[
\langle (M_{t}+i)\tilde f,\tilde u\rangle=0
\qquad
\text{for all }\tilde f\in\mathcal{C}_{\mathrm{res}}^{M},
\]
which says precisely that
\(\tilde u\in((M_{t}+i)\mathcal{C}_{\mathrm{res}}^{M})^{\perp}\).
The converse is the same implication read backwards, so
\eqref{v4-arith:w9-r4:eq:deficiency-range-defect-plus} follows.
The density criterion follows because a subspace has zero orthogonal
complement exactly when its closure is the full Hilbert space.
\end{proof}

\begin{corollary}[range density is graph density for \(M_t\)]
\label{v4-arith:w9-r4:cor:bulk-deficiency-graph-core}
\ClaimStatusProvedHere. For a subspace
\(\mathcal{C}\subset\mathcal{D}(M_t)\), the following are equivalent:
\[
\overline{\mathcal{C}}^{\,\|\cdot\|_{M_t}}=\mathcal{D}(M_t),
\qquad
\overline{(M_t-i)\mathcal{C}}^{\,L^2}=L^2(\mathbb R,dt),
\qquad
\overline{(M_t+i)\mathcal{C}}^{\,L^2}=L^2(\mathbb R,dt).
\]
Consequently, under the residual Meyer chart hypothesis,
\(\mathcal{C}_{\mathrm{res}}^{M}\) is a graph core for \(M_t\) if and
only if
\[
\mathcal{N}_{+}=0=\mathcal{N}_{-}.
\]
\end{corollary}

\begin{proof}
For \(\operatorname{Im}z\neq0\), the map
\[
M_{t}-z\colon \mathcal{D}(M_{t})_{\mathrm{graph}}\to L^{2}(\mathbb{R},dt)
\]
is boundedly invertible: its inverse is multiplication by
\((t-z)^{-1}\), which maps \(L^{2}\) into \(\mathcal{D}(M_{t})\) with
bounded graph norm. Hence a subspace is dense in the graph norm if
and only if its image under \(M_t-z\) is dense in \(L^2\). Apply this
with \(z=i\) and \(z=-i\), then use
\Cref{v4-arith:w9-r4:thm:deficiency-range-defect}.
\end{proof}

\begin{lemma}[rank-one graph-boundary model]
\label{v4-arith:w9-r4:lem:rank-one-graph-boundary}
\ClaimStatusProvedHere. Let
\(\psi(t)=(1+|t|)^{-1}\). Then
\(\psi\in L^{2}(\mathbb{R},dt)\) but
\(\psi\notin\mathcal{D}(M_t)\). Define
\[
\mathcal{C}_{\psi}
:=
\{\,f\in\mathcal{D}(M_t)\mid
\langle (M_t-i)f,\psi\rangle=0\,\}.
\]
Then \(\mathcal{C}_{\psi}\) is \(L^{2}\)-dense and graph-closed of
codimension one, but it is not graph-dense. Its range defects are
\[
\bigl((M_t-i)\mathcal{C}_{\psi}\bigr)^{\perp}
=\mathbb C\psi,
\qquad
\bigl((M_t+i)\mathcal{C}_{\psi}\bigr)^{\perp}
=
\mathbb C\,\frac{t+i}{t-i}\psi .
\]
\end{lemma}

\begin{proof}
The function \(\psi\) is square-integrable, while
\(t\psi(t)\) is not square-integrable. Since
\[
M_t-i\colon\mathcal{D}(M_t)_{\mathrm{graph}}\to L^2(\mathbb R,dt)
\]
is a bounded isomorphism, \(\mathcal{C}_{\psi}\) is the inverse image
of the closed hyperplane \(\psi^\perp\); hence it is graph-closed of
codimension one and its first range defect is \(\mathbb C\psi\).
The defining functional is not continuous for the \(L^2\)-topology on
\(\mathcal{D}(M_t)\), because such continuity would represent it by
an \(L^2\)-vector, equivalently would put \(\psi\) in
\(\mathcal{D}(M_t)\). Therefore its kernel is \(L^2\)-dense.

For the second range, write
\[
(M_t+i)(M_t-i)^{-1}
\]
as multiplication by \((t+i)/(t-i)\), a unitary multiplier on \(L^2\).
The orthogonal complement of its image of \(\psi^\perp\) is the span
of \((t+i)(t-i)^{-1}\psi\).
\end{proof}

\begin{lemma}[point conditions are graph-invisible]
\label{v4-arith:w9-r4:lem:point-conditions-graph-invisible}
\ClaimStatusProvedHere. Let \(\Gamma\subset\mathbb R\) be locally
finite. Then \(C_c^\infty(\mathbb R\setminus\Gamma)\) is graph-dense
in \(\mathcal{D}(M_t)\). In particular, a boundary condition that is
only pointwise at the zero ordinates cannot obstruct graph density for
the ordinary \(M_t\)-graph norm.
\end{lemma}

\begin{proof}
The graph norm of \(M_t\) is the \(L^2\)-norm for the measure
\((1+t^2)\,dt\). Locally finite sets have zero measure for this
weight. Given \(f\in\mathcal{D}(M_t)\), truncate \(f\) to a compact
interval, remove intervals of arbitrarily small weighted measure
around the finitely many points of \(\Gamma\) in that interval, and
then mollify. The approximants lie in
\(C_c^\infty(\mathbb R\setminus\Gamma)\) and converge in the graph
norm.
\end{proof}

\begin{remark}[answer of the calculation]
\label{v4-arith:w9-r4:rem:deficiency-calculation-answer}
The ordinary \(L^{2}\)-bulk part has no hidden deficiency. The only
possible obstruction, once the Tate-polar quotient and residual Meyer
chart are constructed, is boundary-theoretic:
\[
\overline{\mathcal{C}_{\mathrm{res}}^{M}}^{\,\|\cdot\|_{M_{t}}}
\stackrel{?}{=}
\mathcal{D}(M_{t}),
\qquad
\|f\|_{M_{t}}^{2}:=\|f\|_{2}^{2}+\|t f\|_{2}^{2},
\]
in the quotient graph norm. A genuine boundary state must therefore
come from a graph-dual functional of the kind modeled in
\Cref{v4-arith:w9-r4:lem:rank-one-graph-boundary}, not from mere
point conditions at the zero ordinates, which are graph-invisible by
\Cref{v4-arith:w9-r4:lem:point-conditions-graph-invisible}.
\end{remark}

\subsection{Execution Attempt: Graph Norm Versus Determinant}
\label{v4-arith:w9-r4:pdnf:execution-attempt}

We now execute the graph-norm closure request in the currently
explicit ambient spaces. The result is a trichotomy: the classical
Meyer line closes but has the wrong spectrum; the formal zero space
has the right determinant but is tautological; the non-formal
residual quotient is exactly A-TP$^{\mathrm{arch}}$. No non-formal
post-Tate quotient example is presently constructed.

\begin{lemma}[classical Meyer-line graph closure]
\label{v4-arith:w9-r4:lem:meyer-line-graph-closure}
\ClaimStatusProvedHere. Let
\[
M_{t}^{0}:C_{c}^{\infty}(\mathbb{R})\subset L^{2}(\mathbb{R},dt)
\longrightarrow L^{2}(\mathbb{R},dt),
\qquad
(M_{t}^{0}f)(t)=t f(t).
\]
Then \(M_{t}^{0}\) is essentially self-adjoint. Its graph-norm
closure is multiplication by \(t\) on
\[
\mathcal{D}(M_{t})
\;=\;
\{\,f\in L^{2}(\mathbb{R},dt)\mid t f(t)\in L^{2}(\mathbb{R},dt)\,\}.
\]
The spectrum is the absolutely continuous spectrum
\(\mathrm{Spec}(M_{t})=\mathbb{R}\).
\end{lemma}

\begin{proof}
The operator \(M_{t}\) with the displayed maximal domain is
self-adjoint by the spectral theorem for multiplication operators.
The core \(C_{c}^{\infty}(\mathbb{R})\) is dense in the graph norm:
truncate \(f\) by compactly supported cutoffs and then mollify; both
\(f\) and \(t f\) converge in \(L^{2}\). Hence the closure of
\(M_{t}^{0}\) is \(M_{t}\). Since \(M_{t}\) is multiplication by the
coordinate on \(L^{2}(\mathbb{R},dt)\), its spectrum is the essential
range of \(t\), namely \(\mathbb{R}\), with absolutely continuous
spectral measure.
\end{proof}

\begin{proposition}[why the classical closure is insufficient]
\label{v4-arith:w9-r4:prop:classical-closure-insufficient}
\ClaimStatusProvedHere. The closure of
\Cref{v4-arith:w9-r4:lem:meyer-line-graph-closure} does not solve
the spectral-flow problem. Its graph norm is complete and its
operator is self-adjoint, but its spectrum is continuous
\(\mathbb{R}\), not the discrete non-trivial-zero ordinate multiset.
Consequently it cannot satisfy
\[
\det\nolimits_{\infty}
\bigl(s\cdot\mathrm{id}-({\tfrac12}\mathrm{id}+iM_{t})\bigr)
=\xi_{\mathbb{R}}(s)
\]
as a zero-eigenvalue determinant.
\end{proposition}

\begin{proof}
The determinant identity in
\Cref{v4-arith:w9-r4:eq:closure-det} is a regularised product over
the zero-side spectral points. The multiplication operator \(M_{t}\)
on \(L^{2}(\mathbb{R},dt)\) has no such point spectrum; its spectral
measure is Lebesgue measure on the line. Thus the closure obtained
by ordinary Meyer-line graph completion is the unconditional Mellin
unitary part of the story, not the Hadamard-residue spectral
realisation. This is exactly the limitation of Meyer's construction
recorded in \Cref{v4-arith:w9-r4:rem:meyer-schwartz-only}.
\end{proof}

\begin{proposition}[formal zero-space closure is tautological]
\label{v4-arith:w9-r4:prop:formal-zero-closure-tautological}
\ClaimStatusProvedHere. On the formal zero Hilbert space
\(\mathcal{H}_{\xi}^{\mathrm{formal}}\) of
\Cref{v4-arith:w9-r4:def:target-data}, the diagonal ordinate
operator has graph-norm closure and satisfies the regularised
determinant identity if the zero ordinates are used as the input
spectrum. This proves no non-formal spectral-flow closure.
\end{proposition}

\begin{proof}
The finite-support diagonal operator on
\(\bigoplus_{\rho}\mathbb{C}^{m_{\rho}}\) closes in the graph norm to
the maximal diagonal multiplication operator. If all ordinates are
real, the closure is self-adjoint. Given this prescribed spectrum,
\Cref{v4-arith:w5-b:prop:det-formal} identifies the regularised
determinant with \(\xi_{\mathbb{R}}\) by Hadamard factorisation,
up to the stated genus-one normalisation. But the construction has
used the zero multiset as its definition. It supplies no
Bost--Connes, Meyer, Connes, Deninger, or chiral-homological source
for the Hilbert space or the spectral measure; it is the formal
target explicitly excluded by
\Cref{v4-arith:w9-r4:rem:why-formal}.
\end{proof}

\begin{theorem}[execution result for the residual Dirac operator]
\label{v4-arith:w9-r4:thm:execution-result}
\ClaimStatusProvedHere. The requested non-tautological execution on
the residual Meyer/Hadamard core is equivalent to the already named
RH-strength residual. More precisely, the following statements are
equivalent:
\begin{enumerate}
\item The residual operator \(H_{0}=H_{\mathrm{res}}^{\mathrm{sym}}\)
has a graph-norm self-adjoint closure and satisfies
\eqref{v4-arith:w9-r4:eq:closure-det}.
\item The deficiency kernels of
\Cref{v4-arith:w9-r4:def:deficiency-spaces} vanish and the closure
has determinant \(\xi_{\mathbb{R}}\).
\item The Dirac double
\(\mathscr{D}_{\mathrm{res}}^{0}\) admits a Dirac closure in the
sense of \Cref{v4-arith:w9-r4:def:dirac-double}.
\item A-TP$^{\mathrm{arch}}$, equivalently H-P$^{\mathrm{Ar}}$ for
\(\xi_{\mathbb{R}}\).
\end{enumerate}
Thus the execution succeeds on the unquotiented classical line and
on the tautological zero space, but those two successes do not meet
the Vol~IV source-disjointness requirement. On the actual residual
core, the requested proof is the Hilbert--P\'olya/RH-strength step
itself.
\end{theorem}

\begin{proof}
(1) \(\Leftrightarrow\) (2) is von Neumann's deficiency-index
criterion together with the determinant condition, as stated in
\Cref{v4-arith:w9-r4:thm:dirac-two-test}. (2) \(\Leftrightarrow\)
(3) is
\Cref{v4-arith:w9-r4:lem:dirac-double-detects-sa} plus the
definition of Dirac closure. (3) \(\Leftrightarrow\) (4) is the
Polyakov--Dirac normal-form theorem
\Cref{v4-arith:w9-r4:thm:polyakov-dirac-normal-form}. The two
ambient computations above explain why neither classical
Meyer-line closure nor the formal zero Hilbert space supplies the
missing implication: the former has the wrong spectral type, and
the latter imports the desired zero spectrum by definition.
\end{proof}

\begin{remark}[failed proof path recorded]
\label{v4-arith:w9-r4:rem:failed-execution-record}
The graph norm can be closed. The determinant can be obtained. What
cannot be obtained from the present data is both at once on a
non-formal, source-disjoint residual Hilbert space. That conjunction
is precisely the absorption-spectrum property of Connes--Meyer, the
A-TP positivity statement, and the Hilbert--P\'olya realisation.
Any manuscript line claiming this conjunction unconditionally would
be claiming RH.
\end{remark}

\section{Bibliography and Source Catalog}
\label{v4-arith:w9-r4:biblio}

The chapter draws on four catalogs, each internally disjoint and
with explicit disjointness to the catalogs of the upstream Vol~IV
arithmetic chapters.

\paragraph{Hilbert-P\'olya foundational literature.}
\begin{description}
\item[P\'olya.] G.~P\'olya, correspondence with Landau (1914),
transcribed in Pintz and S\'ark\"ozy, \emph{On the work of P\'olya in
analytic number theory}, Stud.~Sci.~Math.~Hungar.~26 (1991), 201--226.
\item[Hilbert.] D.~Hilbert, G\"ottingen seminar notes (1914, 1916),
recorded by E.~Landau, \emph{Handbuch der Lehre von der Verteilung der
Primzahlen}, Teubner (1909), 2nd ed.~1953.
\item[Berry--Keating.] M.~V.~Berry and J.~P.~Keating, \emph{The
Riemann zeros and eigenvalue asymptotics}, SIAM Review 41 (1999),
236--266.
\item[Montgomery.] H.~L.~Montgomery, \emph{The pair correlation of zeros
of the zeta function}, Proc.~Sympos.~Pure Math.~24 (1973), 181--193.
\end{description}

\paragraph{Adelic and idele-class operator constructions.}
\begin{description}
\item[Connes.] A.~Connes, \emph{Trace formula in noncommutative
geometry and the zeros of the Riemann zeta function}, Selecta
Math.~(N.~S.)~5 (1999), 29--106; arXiv:math/9811068. \S IV for the
absorption-spectrum conjecture; \S V for the Weil-positive trace
identity.
\item[Connes--Marcolli.] A.~Connes and M.~Marcolli, \emph{Noncommutative
Geometry, Quantum Fields, and Motives}, AMS Colloq.~Publ.~55 (2008).
Ch.~3 for KMS states; Ch.~4 for rigidity; Ch.~5 for the adele class
space.
\item[Meyer.] R.~Meyer, \emph{On a representation of the idele class
group related to primes and zeros of $L$-functions}, Duke Math.~J.
127 (2005), 519--595; arXiv:math/0311468. \S3 for Mellin transform
and scaling; \S4 for the Schwartz-space distributional trace.
\item[Deninger.] C.~Deninger, \emph{Motivic $L$-functions and
regularized determinants}, Proc.~Sympos.~Pure Math.~55.1 (1994),
707--743; \emph{Some analogies between number theory and dynamical
systems on foliated spaces}, Proc.~ICM~Berlin~1 (1998), 163--186.
\end{description}

\paragraph{Mellin and Hadamard classical analysis.}
\begin{description}
\item[Mellin.] H.~Mellin, \emph{Die Dirichletschen Reihen, die
zahlentheoretischen Funktionen und die unendlichen Produkte von
endlichem Geschlecht}, Acta Soc.~Sci.~Fenn.~31, no.~2 (1904).
\item[Hadamard.] J.~Hadamard, \emph{\'Etude sur les propri\'et\'es
des fonctions enti\`eres et en particulier d'une fonction consid\'er\'ee
par Riemann}, J.~Math.~Pures Appl.~9 (1893), 171--215.
\item[Titchmarsh.] E.~C.~Titchmarsh, \emph{Introduction to the Theory
of Fourier Integrals}, 2nd ed., Oxford (1948); \emph{The Theory of
the Riemann Zeta-Function}, 2nd ed., revised D.~R.~Heath-Brown, Oxford
(1986).
\item[von Mangoldt.] H.~von Mangoldt, \emph{Zu Riemann's Abhandlung
``\"Uber die Anzahl der Primzahlen unter einer gegebenen Gr\"osse''},
J.~Reine Angew.~Math.~114 (1895), 255--305.
\item[Edwards.] H.~M.~Edwards, \emph{Riemann's Zeta Function},
Academic Press (1974). \S2.7 for Hadamard product; \S3 for the
explicit formula.
\item[Patterson.] S.~J.~Patterson, \emph{An Introduction to the Theory
of the Riemann Zeta-Function}, Cambridge (1988). Ch.~3 for the Weil
explicit formula.
\end{description}

\paragraph{Bost--Connes and GNS modular theory.}
\begin{description}
\item[Bost--Connes.] J.-B.~Bost and A.~Connes, \emph{Hecke algebras,
type~III factors and phase transitions with spontaneous symmetry
breaking in number theory}, Selecta Math.~1 (1995), 411--457.
\item[Takesaki.] M.~Takesaki, \emph{Theory of Operator Algebras~II},
Springer (2003). Ch.~VIII for modular automorphism groups.
\item[Julia.] B.~Julia, \emph{Statistical theory of numbers}, Springer
Proc.~Phys.~47 (1990), 276--293. Primon-gas interpretation of the
Bost--Connes spectrum.
\end{description}

\paragraph{First-order and Dirac normal-form sources.}
\begin{description}
\item[Dirac.] P.~A.~M.~Dirac, \emph{The quantum theory of the
electron}, Proc.~Roy.~Soc.~Lond.~A 117 (1928), 610--624, and
118 (1928), 351--361. Source for the first-order square-root
principle used here only as a normal-form discipline.
\item[Reed--Simon.] M.~Reed and B.~Simon, \emph{Methods of Modern
Mathematical Physics I: Functional Analysis}, Academic Press (1980).
Ch.~VIII for adjoints, closures, block operators, and the
self-adjointness criterion used in
\Cref{v4-arith:w9-r4:lem:dirac-double-detects-sa}.
\end{description}

\paragraph{Vol~IV upstream sources used here.}
\begin{description}
\item \Cref{v4-arith:w5-b:hp-operator}: Bost--Connes GNS and
log-modular operator construction (Ch.~\ref{v4-arith:w5-b:hp-operator}).
\item \Cref{v4-arith:rh:hyp:HPAr}: Hilbert--P\'olya hypothesis
(Ch.~\ref{v4-arith:rh-reformulation}).
\item \Cref{v4-arith:HPAr-vbconverse}: wave-2 equivalence bundle
(Ch.~\ref{v4-arith:HPAr-vbconverse}).
\item \Cref{v4-arith:hpar-det}: wave-3 determinacy ledger
(Ch.~\ref{v4-arith:hpar-det}).
\item \Cref{v4-arith:w7-j:ATP-adelic-scattering}: adelic-scattering
formulation (Ch.~\ref{v4-arith:w7-j:ATP-adelic-scattering}).
\item \Cref{v4-arith:w8-a:deninger-chiral-homology}: Deninger-chiral
triple identification (Ch.~\ref{v4-arith:w8-a:deninger-chiral-homology}).
\item \Cref{v4-arith:3d-acs-E1bar}: 3D arithmetic Chern-Simons
(Ch.~\ref{v4-arith:3d-acs-E1bar}).
\end{description}

\paragraph{Disjointness audit.}
The present chapter draws its formal residue-obstruction analysis
(\Cref{v4-arith:w9-r4:lem:formal-intertwine}) from the
Riemann--Hadamard--von Mangoldt classical catalog, disjoint from the
Bost--Connes / Takesaki modular-theoretic catalog of
Ch.~\ref{v4-arith:w5-b:hp-operator}; its Meyer spectral-matching
analysis from Meyer 2005 arXiv:math/0311468, disjoint from the
Birman--Krein scattering catalog of
Ch.~\ref{v4-arith:w7-j:ATP-adelic-scattering}; its Connes absorption
analysis from Connes 1999 \S IV, disjoint from the Bost--Connes 1995
$C^{*}$-algebraic catalog of Ch.~\ref{v4-arith:w5-b:hp-operator}. No
derivation on which a residual-irreducibility claim rests is shared
between the catalogs.
